% appendix_convergence.tex: fragment for \appendix% appendix_convergence.tex: fragment for \appendix% appendix_convergence.tex: fragment for \appendix% appendix_convergence.tex: fragment for \appendix\input{appendix_convergence} in note.tex.
% Packages: amsmath and amssymb only (already loaded by note.tex); no amsthm.
% Labels of note.tex used: setting, estimator, tests, limits, clock, cov, model, square, ball.
% Citation keys used: brewin, gibbons, wang, madsen, and ONE new key, oneill:
%   \bibitem{oneill} O'Neill, B.: Semi-Riemannian Geometry with Applications to
%   Relativity. Academic Press, New York (1983)
% Check of the orders and of the rank certificate: python check_convergence_orders.py
\section{A conditional convergence statement}\label{app:conv}
This appendix states what we can prove about the fit of
Sect.~\ref{estimator} in shrinking windows, and what we cannot. The
statement is conditional: it assumes that the landmarks form a
well-conditioned configuration, and it concerns the solution of the fit that
lies near the true geometry. It does not show that the selection rule and
the optimizer of Sect.~\ref{estimator} meet these two conditions, and it
says nothing about the coverage of the ball \eqref{ball}. Only orders are
established; no constant is evaluated.

\subsection{Setting and assumptions}\label{app:setting}
We write $|\cdot|$ for the Euclidean norm of a vector of components,
$\|\cdot\|$ for the corresponding operator norm, and $A\succeq B$ if $A-B$
is positive semidefinite. For two events $x,x'$ of a convex normal
neighbourhood, $s(x,x')=g(v,v)$ with $v=\exp_x^{-1}x'$ is the signed squared
geodesic interval, so that $s=-\tau^2$ for timelike separation, and the
quantity $\sigma(x,y)$ of Sect.~\ref{estimator} is $s/T^2$.

\begin{description}
\item[(A1) Geometry.] $(M,g)$ is a time-oriented four-dimensional Lorentzian
manifold with smooth metric, $o\in M$, and $(e_a)$ is an orthonormal frame
at $o$ with $e_0$ future-directed. $U$ is a convex normal neighbourhood of
$o$, that is, an open set that is a normal neighbourhood of each of its
points; then any two points of $U$ are joined by a unique geodesic segment
in $U$, and $(x,x')\mapsto\exp_x^{-1}x'$ is smooth on $U\times U$
\cite[Chap.~5]{oneill}. The metric is Ricci-flat on $U$, and in the normal
coordinates based at $(o,e_a)$ its matrix of components satisfies
$\|g_{ab}(x)-\eta_{ab}\|\le1/10$ on $U$. Every point of a Ricci-flat region
has arbitrarily small neighbourhoods of this kind. For $0<T\le T_0$ the tips
are $p=\exp_o(-\tfrac T2e_0)$ and $q=\exp_o(\tfrac T2e_0)$, and, as in
Sect.~\ref{setting}, $I(p,q)=J^+(p)\cap J^-(q)\subset U$; this holds for
small $T$ if $M$ is strongly causal at $o$. The proofs use finitely many
derivatives, six of $s$ and four of the metric components in normal
coordinates; we have not determined the least differentiability of $g$ that
provides them.
\item[(A2) Sampling.] Conditionally on the landmarks, the $N$ measurement
events are independent and uniformly distributed on $I(p,q)$ with respect to
the volume of $g$. This is the case when the landmarks are taken from an
independent selection sample.
\item[(A3) Landmark configuration.] There are $L$ landmarks
$\ell_1,\dots,\ell_L\in I(p,q)$; put $\ell_0=p$ and $\ell_{L+1}=q$. Their
normal coordinates in units of $T$ are $x_i\in\mathbb R^4$, that is,
$\ell_i=\exp_o(Tx_i^ae_a)$, so that $x_0=(-\tfrac12,0,0,0)$ and
$x_{L+1}=(\tfrac12,0,0,0)$. The spatial axes are those of the gauge of
Sect.~\ref{estimator}: $x_1=(x_1^0,x_1^1,0,0)$ with $x_1^1>0$ and
$x_2=(x_2^0,x_2^1,x_2^2,0)$ with $x_2^2>0$. This fixes the frame up to the
reflection of $e_3$, and we fix one of the two frames. The remaining $4L-3$
coordinates form $\xi\in\mathbb R^{4L-3}$. The clocks are the pairs
$(\ell_i,\ell_j)$ with $(i,j)$ in a fixed set $\mathcal P$ of $m$ index
pairs, $(i,j)\neq(0,L+1)$. We assume that $\xi$ lies in a fixed compact set
$\mathcal K\subset\mathbb R^{4L-3}$ such that, for every configuration in
$\mathcal K$, the $L+2$ points $x_i$ are distinct, $x_1,\dots,x_L$ lie in
the open flat interval $\{|x^0|+|\vec x|<\tfrac12\}$, and $x_j-x_i$ is
future-directed timelike with respect to $\eta$ for every
$(i,j)\in\mathcal P$.
\item[(A4) Conditioning.] Let $\theta=(\xi,c,d)\in\mathbb R^D$, $D=4L+31$,
where $c\in\mathbb R^{10}$ and $d\in\mathbb R^{24}$ are the components of the
curvature and of its gradient in the bases of Sect.~\ref{estimator} (there
$d$ is called $g$, a letter needed here for the metric). Let
$\sigma(\theta)\in\mathbb R^m$ be the model \eqref{model} for the pairs of
$\mathcal P$, with $\sigma_2$ as in \eqref{square} or with $\sigma_2=0$, and
$J(\theta)=\partial\sigma/\partial\theta$. We assume that the $m\times D$
matrix $J(\xi,0,0)$ has rank $D$ for every $\xi\in\mathcal K$.
\item[(A5) Weights.] The fit minimizes
$\chi^2(\theta)=N[y-\sigma(\theta)]^{\!\top}\widehat W[y-\sigma(\theta)]$
with a symmetric matrix $\widehat W$ that may depend on the data, and there
are constants $0<w_-\le w_+$ such that $w_-I\preceq\widehat W\preceq w_+I$
with probability at least $1-\varpi_N$, where $\varpi_N\to0$ as
$N\to\infty$. For the weight of Sect.~\ref{estimator},
$\widehat W=(N\widehat\Sigma)^{-1}$, this is proved in
Lemma~4 below.
\end{description}

Assumptions (A3) and (A4) are substantive, and (A4) is a condition on the
shape of the configuration only. By compactness there are $\tau_0>0$ and
$\mu>0$ such that $\eta(x_j-x_i,x_j-x_i)\le-\tau_0^2$ for
$(i,j)\in\mathcal P$ and the smallest singular value of $J(\xi,0,0)$ is at
least $\mu$, for all $\xi\in\mathcal K$; all constants below depend on
these numbers. The spatial axes of the frame $(e_a)$ depend on the
landmarks; since the orthonormal frames at $o$ with the given $e_0$ form a
compact set, all estimates below are uniform over them. The true parameter
is $\theta^*_T=(\xi,c^*,d^*)$, where $\xi$
is the true configuration and $c^*$, $d^*$ are the components of
$T^2C_{abcd}(o)$ and of $T^3\nabla_eC_{abcd}(o)$ in the frame $(e_a)$. The
measured clocks are $y_{ij}=-(n_{ij}/N)^{1/2}$, with $n_{ij}$ the number of
measurement events in $I(\ell_i,\ell_j)$, and the continuum clocks are
$\bar y_{ij}=-[V(\ell_i,\ell_j)/V(p,q)]^{1/2}$.

\subsection{Statement}
\noindent\textbf{Proposition 1.} \emph{Assume (A1)--(A5). There are
constants $T_1,r_1>0$ and $C_1,C_2<\infty$, which depend only on the
geometry of $(U,g)$, on $\mathcal K$, $\mathcal P$ and on $w_\pm$, such that
for all $0<T\le T_1$, all $N\ge2$ and every landmark configuration allowed
by (A3) the following holds with probability at least $1-\varpi_N-C_2/N$.}
\begin{enumerate}
\item[(i)] \emph{$\chi^2$ has exactly one critical point
$\hat\theta=(\hat\xi,\hat c,\hat d)$ in the closed ball of radius $r_1$
around $\theta^*_T$, and $\hat\theta$ is the unique minimizer of $\chi^2$
on this ball.}
\item[(ii)] \emph{$|\hat\theta-\theta^*_T|\le C_1\left(T^4+|y-\bar y|\right)$.}
\end{enumerate}
\emph{Moreover $\operatorname{E}|y-\bar y|^2\le m/N$. Consequently, for
every $t>0$, with probability at least $1-\varpi_N-C_2/N-m/t^2$,}
\begin{equation}
 |\hat\theta-\theta^*_T|\le C_1\left(T^4+tN^{-1/2}\right),\qquad
 \left|\hat c/T^2-c^*/T^2\right|\le C_1\left(T^2+tN^{-1/2}T^{-2}\right),
 \label{conv:rate}
\end{equation}
\emph{where $c^*/T^2$ are the components of the Weyl tensor $C_{abcd}(o)$.
The same holds, with other constants, in the tidal norm of
Sect.~\ref{setting}.}

\medskip
In short, $\hat\theta-\theta^*_T=O(T^4)+O_p(N^{-1/2})$ in the normalized
variables of the fit. The statement and its proof are the same with and
without $\sigma_2$, because $\sigma_2=O(|c|^2)$; in particular the
coefficient of the volume term in \eqref{square} is not used. The estimator
returns $\hat c$, an estimate of $T^2C(o)$. The second inequality of
\eqref{conv:rate} refers to the tensor in physical units, which requires
$T$, or equivalently bounds the relative error when $C(o)\neq0$.

\medskip
\noindent\textbf{Corollary 1.} \emph{Let the measurement events be a
Poisson sprinkling of density $\rho$ in $I(p,q)$, with lengths measured in
a fixed unit, and let (A1)--(A5) hold. If $T\to0$ and $\rho T^8\to\infty$,
then the solution $\hat\theta$ of Proposition~1 exists with probability
tending to one, and}
\begin{equation}
 \hat c/T^2-C(o)=O(T^2)+O_p\!\left(\rho^{-1/2}T^{-4}\right)\to0
 \quad\text{in probability}.
\end{equation}
\emph{For $T$ proportional to $\rho^{-1/12}$ the error is
$O_p(\rho^{-1/6})$.}

\subsection{Proofs}
\noindent\textbf{Lemma 1} (normal coordinates). \emph{Let $\mathcal F$ be a
compact set of orthonormal frames $(x,f)$ with $x\in U$. Let
$w\in\mathbb R^4$ be the normal coordinates based at $(x,f)$, and let
$s_{x,f}(w,w')$ and $\gamma_{x,f}(w)$ be the function $s$ and the matrix of
metric components in these coordinates. There are $r_0>0$ and $C_0$ such
that for all $(x,f)\in\mathcal F$ and $|w|,|w'|\le r_0$, with
$\Delta=w-w'$:}
\begin{enumerate}
\item[(a)] \emph{$s_{x,f}(w,w')=\eta(\Delta,\Delta)-\tfrac13R(w,w',w,w')
-\tfrac1{12}(w+w')^e\nabla_eR(w,w',w,w')+\varrho$ with
$|\varrho|\le C_0(|w|+|w'|)^6$, where $R$ and $\nabla R$ are the components
of the curvature and its gradient at $x$ in the frame $f$;}
\item[(b)] \emph{$s_{x,f}(w,w')-\eta(\Delta,\Delta)$ is bounded by
$C_0(|w|+|w'|)^4$ and its first derivatives by $C_0(|w|+|w'|)^3$;}
\item[(c)] $\bigl|\,|\det\gamma_{x,f}(w)|^{1/2}-1\,\bigr|\le C_0|w|^4$;
\item[(d)] $s_{x,f}(0,w)=\eta(w,w)$.
\end{enumerate}

\noindent\emph{Proof.} The maps $(x,f,w,w')\mapsto s_{x,f}(w,w')$ and
$(x,f,w)\mapsto\gamma_{x,f}(w)$ are smooth on an open set that contains
$\mathcal F\times\{0\}\times\{0\}$, hence on
$\mathcal F\times\{|w|,|w'|\le r_0\}$ for some $r_0>0$. Statement (d) is the
definition of $s$. For (a) apply Taylor's theorem in $(w,w')$ at the
origin: the derivatives of order six are bounded on this compact set, and
the Taylor polynomial of degree five is the expansion of the squared
geodesic interval in normal coordinates, Sect.~9 of Ref.~\cite{brewin}
(Eqs.~(11.21a)--(11.21c) in arXiv:0903.2087v3). That polynomial has no terms
of degree three, which gives (b) in the same way. For (c),
$\gamma_{ab}(w)=\eta_{ab}-\tfrac13R_{acbd}w^cw^d
-\tfrac16\nabla_eR_{acbd}w^cw^dw^e+O(|w|^4)$ \cite{brewin}, and the traces
of the second and third terms with $\eta^{ab}$ contain only the Ricci tensor
and its gradient, which vanish; hence $\det\gamma=-1+O(|w|^4)$, with a
remainder bounded as before. \hfill$\square$

\medskip
\noindent\textbf{Lemma 2} (volume of a small vacuum interval). \emph{Let
$\mathcal T$ be a compact set of future-directed unit timelike vectors
$(x,u)$ with $x\in U$. There are $\tau_1>0$ and $C_V$ such that, whenever
$(x,u)\in\mathcal T$, $0<\tau\le\tau_1$, $x'=\exp_x(\tau u)$ and
$J^+(x)\cap J^-(x')\subset U$,}
\begin{equation}
 \left|V(x,x')-\tfrac{\pi}{24}\tau^4\right|\le C_V\tau^8 .
 \label{conv:volume}
\end{equation}

\noindent\emph{Proof.} Complete $u$ to an orthonormal frame $f$ with
$f_0=u$. The frames so obtained form a compact set, to which Lemma~1
applies. We use the normal coordinates $w$ based at $(x,f)$, in which $x'$
is $\tau\hat e$ with $\hat e=(1,0,0,0)$, write $s=s_{x,f}$, and put
$\Xi(w)=s(0,w)-s(w,\tau\hat e)$. We show that the half
$H_-=\{z\in J^+(x)\cap J^-(x'):\Xi\ge0\}$ of the interval has volume
$\tfrac{\pi}{48}\tau^4+O(\tau^8)$. The other half, where $\Xi\le0$, is
treated in the same way from $x'$ with the time orientation reversed, and
the common boundary $\Xi=0$ has measure zero. All $O(\cdot)$ below are
uniform over $\mathcal T$ and over the directions $\nu$.

(1) \emph{Localization.} By (A1), every causal vector $v$ at a point of $U$
has $|\vec v|\le\tfrac54|v^0|$ in the coordinates of (A1), with $v^0>0$ if
it is future-directed. Hence $x^0$ increases along future-directed causal
curves in $U$, and the coordinate differences satisfy
$|z-x|\le\tfrac94|x'-x|$ for every $z\in J^+(x)\cap J^-(x')$. Since the
change to the coordinates $w$ is smooth, uniformly over the frames, there
is $C\ge1$ such that $|w|\le C\tau$ for all these $z$.

(2) \emph{Rays.} In a convex normal neighbourhood, $z\in J^+(x)$ if and
only if $\exp_x^{-1}z$ is zero or future-directed causal
\cite[Lemma~14.2]{oneill}, that is, if and only if $w=t\nu$ with $t\ge0$
and $\nu\in\mathcal N=\{(1,\vec\nu):|\vec\nu|\le1\}$. For $\nu\in\mathcal N$
let $\phi_\nu(t)=\Xi(t\nu)$. By Lemma~1(d),
$\phi_\nu(t)=t^2\eta(\nu,\nu)-s(t\nu,\tau\hat e)$, and
$\eta(t\nu-\tau\hat e,t\nu-\tau\hat e)=t^2\eta(\nu,\nu)+2\tau t-\tau^2$.
By Lemma~1(b), for $0\le t\le C\tau$,
\begin{equation}
 \phi_\nu(t)=\tau^2-2\tau t+O(\tau^4),\qquad
 \phi_\nu'(t)=-2\tau+O(\tau^3).
\end{equation}
For $\tau\le\tau_1$ small, $\phi_\nu$ is therefore strictly decreasing on
$[0,C\tau]$, with $\phi_\nu(0)=\tau^2$ and $\phi_\nu(\tau)<0$. It has a
single zero $t_*(\nu)\in(0,\tau)$, which depends continuously on $\nu$.

(3) \emph{The half-interval is
$\{t\nu:\nu\in\mathcal N,\ 0\le t\le t_*(\nu)\}$.} If $z\in H_-$, then
$w=t\nu$ with $t\le C\tau$ by (1) and (2), and $\phi_\nu(t)\ge0$ gives
$t\le t_*(\nu)$. Conversely, let $0\le t\le t_*(\nu)$. On the segment from
$0$ to $t\nu$ we have $\Xi\ge0$, that is, $s(\cdot,x')\le s(x,\cdot)\le0$,
so the segment lies in the set of events causally related to $x'$, and it
does not contain $x'$, where $\Xi=-\tau^2$. The sets
$J^-(x')\setminus\{x'\}$ and $J^+(x')\setminus\{x'\}$ are disjoint and
closed in $U\setminus\{x'\}$, and the segment starts at $x\in J^-(x')$;
being connected, it lies in $J^-(x')$. Hence $t\nu\in H_-$.

(4) \emph{Volume.} With $w=(t,t\vec\nu)$ we have
$\mathrm d^4w=t^3\,\mathrm dt\,\mathrm d^3\vec\nu$, and Lemma~1(c) gives
\begin{equation}
 \operatorname{Vol}(H_-)=\int_{|\vec\nu|\le1}\int_0^{t_*(\nu)}
 \left[1+O(\tau^4)\right]t^3\,\mathrm dt\,\mathrm d^3\vec\nu
 =\frac14\int_{|\vec\nu|\le1}t_*(\nu)^4\,\mathrm d^3\vec\nu+O(\tau^8).
 \label{conv:half}
\end{equation}

(5) \emph{Expansion of $t_*$.} Insert Lemma~1(a) with $w=t\nu$ and
$w'=\tau\hat e$. Since
$R(t\nu,\tau\hat e,t\nu,\tau\hat e)=t^2\tau^2E_{ij}\nu^i\nu^j$ with
$E_{ij}=R_{0i0j}$,
\begin{equation}
 \phi_\nu(t)=\tau^2-2\tau t+\tfrac13t^2\tau^2E_{ij}\nu^i\nu^j
 +\tfrac1{12}t^2\tau^2\left[(t+\tau)\nabla_0E_{ij}
 +t\,\nu^k\nabla_kE_{ij}\right]\nu^i\nu^j+O(\tau^6).
\end{equation}
The zero satisfies $t_*=\tfrac\tau2+O(\tau^3)$, and then
\begin{equation}
 t_*(\nu)=\frac\tau2\left[1+\frac{\tau^2}{12}E_{ij}\nu^i\nu^j
 +\frac{\tau^3}{96}\left(3\nabla_0E_{ij}+\nu^k\nabla_kE_{ij}\right)
 \nu^i\nu^j+O(\tau^4)\right].
\end{equation}
The integral of $\nu^i\nu^j$ over the unit ball is
$\tfrac{4\pi}{15}\delta^{ij}$, and that of $\nu^i\nu^j\nu^k$ vanishes. In
vacuum $E_{ii}=R_{00}=0$ and $\nabla_0E_{ii}=\nabla_0R_{00}=0$, so the terms
of relative order $\tau^2$ and $\tau^3$ integrate to zero, and
\eqref{conv:half} gives
$\operatorname{Vol}(H_-)=\tfrac14\cdot\tfrac{4\pi}3(\tau/2)^4+O(\tau^8)
=\tfrac{\pi}{48}\tau^4+O(\tau^8)$. \hfill$\square$

\medskip
Lemma~2 uses the field equations only through $R_{ab}=0$, and it gives only
the order of the correction. As a check of the method, without that
condition the same computation, with the volume element
$1-\tfrac16R_{ab}w^aw^b$, gives
$V=\tfrac{\pi}{24}\tau^4[1-\tfrac1{180}R^a{}_a\tau^2+\tfrac1{30}R_{00}\tau^2
+O(\tau^3)]$, which is the formula of Myrheim, Eqs.~(7), (14) and (23) of
Ref.~\cite{gibbons}. In vacuum the first correction is of relative order
$\tau^4$, in agreement with Eq.~(79) of Ref.~\cite{wang}; its coefficient, on
which \eqref{square} and Ref.~\cite{wang} differ, plays no role here.

\medskip
\noindent\textbf{Lemma 3} (clocks). \emph{Assume (A1)--(A3). (a) There are
$T_1>0$ and $C_b$ such that $|\bar y_{ij}-\sigma_{ij}(\theta^*_T)|\le
C_bT^4$ for all $T\le T_1$, $\xi\in\mathcal K$ and $(i,j)\in\mathcal P$.
(b) $\operatorname{E}|y-\bar y|^2\le m/N$.}

\noindent\emph{Proof.} (a) Lemma~1(a) at $(o,e)$ with $w=Tx_i$ and
$w'=Tx_j$, and $R_{abcd}=C_{abcd}$ in vacuum, give
\begin{equation}
 \frac{s(\ell_i,\ell_j)}{T^2}=\eta(\Delta,\Delta)
 -\frac{T^2}3C(x_i,x_j,x_i,x_j)
 -\frac{T^3}{12}(x_i+x_j)^e\nabla_eC(x_i,x_j,x_i,x_j)+O(T^4),
\end{equation}
with $\Delta=x_i-x_j$ and the tensors at $o$. The right-hand side is the
model \eqref{model} at $\theta^*_T$ without $\sigma_2$, up to $O(T^4)$, and
$\sigma_2(\theta^*_T)=O(|c^*|^2)=O(T^4)$. Since
$\eta(\Delta,\Delta)\le-\tau_0^2$, for small $T$ we have
$\tau_{ij}^2=-s(\ell_i,\ell_j)\ge\tau_0^2T^2/2$ and $\tau_{ij}=O(T)$. In
the coordinates of (A1) the vector $\exp_{\ell_i}^{-1}\ell_j$ has components
$T(x_j-x_i)+O(T^2)$, so for small $T$ it is future-directed and $\ell_j$
lies in the timelike future of $\ell_i$. The unit tangent at $\ell_i$ of the
geodesic to $\ell_j$ then has bounded components, and these vectors lie in a
compact set $\mathcal T$.
Moreover $J^+(\ell_i)\cap J^-(\ell_j)\subset I(p,q)\subset U$. Lemma~2
therefore gives $V(\ell_i,\ell_j)=\tfrac{\pi}{24}\tau_{ij}^4[1+O(T^4)]$, and
likewise $V(p,q)=\tfrac{\pi}{24}T^4[1+O(T^4)]$. Hence
$\bar y_{ij}=-(\tau_{ij}^2/T^2)[1+O(T^4)]=s(\ell_i,\ell_j)/T^2+O(T^4)$.

(b) Given the landmarks, $n_{ij}$ is binomial with parameters $N$ and
$\pi_{ij}=V(\ell_i,\ell_j)/V(p,q)=\bar y_{ij}^2$. Since
$|\sqrt a-\sqrt b|\le|a-b|/\sqrt b$,
\begin{equation}
 |y_{ij}-\bar y_{ij}|\le\frac{|n_{ij}/N-\pi_{ij}|}{\sqrt{\pi_{ij}}},
 \qquad\operatorname{E}|y_{ij}-\bar y_{ij}|^2\le\frac{1-\pi_{ij}}N .
\end{equation}
Summing over the $m$ pairs gives (b). The clocks are strongly correlated,
see \eqref{cov}, but no independence is used. \hfill$\square$

\medskip
\noindent\textbf{Lemma 4} (weights). \emph{Let $N\widehat\Sigma=\widehat
D\,\widehat\Gamma\,\widehat D$ be the covariance of Sect.~\ref{estimator},
where $\widehat\Gamma$ is the empirical covariance matrix of the membership
indicators of the $m$ intervals and
$\widehat D=\operatorname{diag}(\tfrac12\hat\pi_{ij}^{-1/2})$ with
$\hat\pi_{ij}=n_{ij}/N$. Under (A1)--(A3) there are $T_2>0$,
$0<w_-\le w_+$ and $C_W$ such that $\widehat W=(N\widehat\Sigma)^{-1}$
satisfies (A5) with $\varpi_N=C_W/N$ for all $T\le T_2$.}

\noindent\emph{Proof.} Let $\Gamma_T(\xi)$ be the covariance matrix of the
indicators under the sampling law, and $\Gamma_0(\xi)$ the same matrix for
the flat metric $\eta$ and the points $x_i$.

(1) \emph{$\Gamma_0(\xi)$ is positive definite.} Otherwise a combination
$\sum\lambda_{ij}1_{ij}$ of the indicators with $\lambda\neq0$ is constant
almost everywhere in the flat interval. Fix a pair $(i,j)$. The edge of its
interval, where the future cone of $x_i$ meets the past cone of $x_j$, is a
round two-sphere, and the only null cones that contain it are those of
$x_i$ and $x_j$; the null cone of any other point meets it in the zero set
of a quadratic polynomial that does not vanish identically on the sphere,
a closed set without interior points. Choose $z$ on the edge and on none of
the other cones, and a ball $B$ around $z$ that meets no null cone of a
point $x_k$ other than the future cone of $x_i$ and the past cone of $x_j$.
In $B$, every indicator other than $1_{ij}$ is a constant, or a constant
times the indicator of $J^+(x_i)$ or of $J^-(x_j)$, while $1_{ij}$ is the
product of these two. The two cones cross transversally at $z$. If
$1\le i,j\le L$, all four combinations of the two indicators occur on sets
of positive measure in $B$, and constancy of the sum forces
$\lambda_{ij}=0$. If $i=0$, only $B\cap J^+(p)$ carries probability, and no
interval $I(x_k,x_j)$ with $k\neq0$ meets $B$, because
$z\in\partial J^+(p)$ whereas $x_k$ lies in the open interval; the two
values of the indicator of $J^-(x_j)$ again force $\lambda_{0j}=0$. The
case $j=L+1$ is symmetric.

(2) In the coordinates $x/T$ the metric has components $g_{ab}(Tx)$. These,
the inverse exponential map and the volume density converge to their flat
values as $T\to0$, uniformly on compact sets. The indicator of each
interval is therefore continuous in $(T,\xi)$ and in the spatial frame at
$o$ at every point off the boundary of the limiting interval, and by
dominated convergence $\Gamma_T(\xi)$ is continuous in these variables on
$[0,T_0]\times\mathcal K$ times the compact set of frames, with the value
$\Gamma_0(\xi)$ at $T=0$. With (1) and compactness,
$\Gamma_T(\xi)\succeq4\gamma I$ for some $\gamma>0$, all $\xi\in\mathcal K$,
all frames and all $T\le T_2$. In the same way $\pi_{ij}\ge\pi_->0$.

(3) $\widehat\Gamma$ and $\hat\pi$ are built from averages of bounded
variables, so $\operatorname{E}\|\widehat\Gamma-\Gamma_T\|^2\le Cm^2/N$ and
$\operatorname{E}|\hat\pi-\pi|^2\le m/N$. By Chebyshev's inequality the
event that $\|\widehat\Gamma-\Gamma_T\|\le2\gamma$ and
$\hat\pi_{ij}\ge\pi_-/2$ for all pairs has probability at least $1-C_W/N$.
On this event the diagonal entries of $\widehat D$ lie between $\tfrac12$
and $(2\pi_-)^{-1/2}$, $\widehat\Gamma\succeq2\gamma I$ and
$\operatorname{tr}\widehat\Gamma\le m/2$, hence
$\tfrac\gamma2I\preceq N\widehat\Sigma\preceq\tfrac{m}{4\pi_-}I$.
\hfill$\square$

\medskip
\noindent\textbf{Lemma 5} (local minimizer). \emph{Let $F$ be a twice
continuously differentiable map from the closed ball of radius $r_0$ around
$\theta^*\in\mathbb R^D$ to $\mathbb R^m$, with $\|DF\|\le M_1$ and
$\|D^2F\|\le M_2$. Let $W$ be symmetric with $0\preceq W\preceq w_+I$ and
$DF(\theta^*)^{\!\top}W\,DF(\theta^*)\succeq\lambda I$, $\lambda>0$. For
$y\in\mathbb R^m$ let $Q(\theta)=\tfrac12[y-F(\theta)]^{\!\top}
W[y-F(\theta)]$ and $\epsilon=|y-F(\theta^*)|$, and put}
\begin{equation}
 r_1=\min\left\{r_0,\frac{\lambda}{12M_1M_2w_+}\right\},\qquad
 \epsilon_1=\min\left\{\frac{\lambda}{4M_2w_+},
 \frac{\lambda r_1}{8M_1w_+}\right\}.
\end{equation}
\emph{If $\epsilon\le\epsilon_1$, then $Q$ has exactly one critical point
$\hat\theta$ in the closed ball of radius $r_1$ around $\theta^*$, it is the
unique minimizer of $Q$ on this ball, and
$|\hat\theta-\theta^*|\le4M_1w_+\epsilon/\lambda$.}

\noindent\emph{Proof.} Let $|\theta-\theta^*|\le r_1$, let $v$ be a unit
vector and $\rho=y-F(\theta)$. Then
\begin{equation}
 v^{\!\top}D^2Q(\theta)\,v=[DF(\theta)v]^{\!\top}W[DF(\theta)v]
 -\rho^{\!\top}W\,D^2F(\theta)[v,v].
\end{equation}
Here $|\rho|\le\epsilon+M_1r_1$ and $\|DF(\theta)-DF(\theta^*)\|\le M_2r_1$,
so that
$v^{\!\top}D^2Q(\theta)v\ge\lambda-2M_1M_2w_+r_1-M_2w_+(\epsilon+M_1r_1)
\ge\lambda/2$. Thus $Q$ is strongly convex on the ball. It has a unique
minimizer $\hat\theta$ there and at most one critical point, and
$Q(\theta)\ge Q(\theta^*)+DQ(\theta^*)(\theta-\theta^*)
+\tfrac\lambda4|\theta-\theta^*|^2$. Since
$|DQ(\theta^*)|=|DF(\theta^*)^{\!\top}W[y-F(\theta^*)]|\le M_1w_+\epsilon$,
we get $Q(\theta)>Q(\theta^*)$ whenever
$|\theta-\theta^*|>4M_1w_+\epsilon/\lambda$. Hence
$|\hat\theta-\theta^*|\le4M_1w_+\epsilon/\lambda\le r_1/2$, so $\hat\theta$
is an interior point and therefore a critical point. \hfill$\square$

\medskip
\noindent\emph{Proof of Proposition 1.} The map $\sigma$ is a polynomial in
$\theta$. Let $M_1$ and $M_2$ bound its first and second derivatives on the
set of all $\theta$ at distance at most one from
$\{(\xi,c,d):\xi\in\mathcal K,\ |c|,|d|\le1\}$; for small $T$ this set
contains the ball of radius $r_0=1$ around $\theta^*_T$. Because $\sigma_2$
is quadratic in $c$, the matrix $J(\xi,0,0)$ does not depend on it, and
$\|J(\theta^*_T)-J(\xi,0,0)\|\le M_2(|c^*|+|d^*|)=O(T^2)$. By (A4) the
smallest singular value of $J(\theta^*_T)$ is at least $\mu/2$ for
$T\le T_1$, so that on the event of (A5)
$J(\theta^*_T)^{\!\top}\widehat WJ(\theta^*_T)\succeq\lambda I$ with
$\lambda=w_-\mu^2/4$. Apply Lemma~5 on this event with $F=\sigma$,
$W=\widehat W$ and $\theta^*=\theta^*_T$; the functions $\chi^2$ and $Q$
differ by the factor $2N$. By Lemma~3(a),
$\epsilon=|y-\sigma(\theta^*_T)|\le|y-\bar y|+\sqrt m\,C_bT^4$. Reduce $T_1$
so that $\sqrt m\,C_bT_1^4\le\epsilon_1/2$. By Lemma~3(b) and Chebyshev's
inequality, $|y-\bar y|\le\epsilon_1/2$ except on an event of probability at
most $C_2/N$ with $C_2=4m/\epsilon_1^2$. Lemma~5 then gives (i), and (ii)
with $C_1=4M_1w_+\max(1,\sqrt m\,C_b)/\lambda$. Finally
$|y-\bar y|\le tN^{-1/2}$ except with probability $m/t^2$, which gives
\eqref{conv:rate}, the second inequality because
$|\hat c-c^*|\le|\hat\theta-\theta^*_T|$. \hfill$\square$

\medskip
\noindent\emph{Proof of Corollary 1.} Given their number $N$, the events of
a Poisson sprinkling in $I(p,q)$ satisfy (A2), and $N$ is Poisson with mean
$\rho V(p,q)=\tfrac{\pi}{24}\rho T^4[1+O(T^4)]$ by Lemma~2. If
$\rho T^4\to\infty$, then $N\ge\rho T^4/20$ with probability tending to
one, and on this event $N^{-1/2}T^{-2}\le\sqrt{20}\,\rho^{-1/2}T^{-4}$ and
$\varpi_N+C_2/N\to0$. Insert this into \eqref{conv:rate}. The two terms
balance for $T^6$ proportional to $\rho^{-1/2}$, and then
$\rho T^8\to\infty$. \hfill$\square$

\subsection{What is not proved}\label{app:open}
\emph{The selection rule.} Assumptions (A3) and (A4) are not shown to hold
for the landmarks produced by the rule of Sect.~\ref{estimator}, which are
events of a selection sample chosen by counts, with a clock set chosen by
counts. Assumption (A4) is not vacuous: for an explicit rational
configuration with five landmarks in each of three time bands, and with
the clocks between different bands and to the tips, exact modular
arithmetic gives $J(\xi,0,0)$ the full rank $D=91$. Since full rank is the
non-vanishing of a polynomial, (A4) then holds for almost every
configuration with these clocks (see Corollary~2). This identifies the
parameters; it does not bound $\mu$, on which the constants depend. In
that configuration the singular values of $J(\xi,0,0)$ range over eight
orders of magnitude.

\emph{The optimizer.} Proposition~1 concerns the critical point near
$\theta^*_T$. It does not show that the Levenberg--Marquardt iteration from
the four starts of Sect.~\ref{estimator} converges to this point, and it
does not exclude other local minima. The reflections of
Sect.~\ref{setting} map $\hat\theta$ to further minimizers with the same
$\chi^2$; these are the mirror images accounted for there. No global
identifiability is proved.

\emph{The coverage.} Nothing is proved about the ball \eqref{ball}. For
$T$ proportional to $\rho^{-1/12}$ the bias and the noise in
\eqref{conv:rate} are of the same order, so a ball built from the noise
covariance alone cannot be expected to cover, and the allowance $\delta$
remains an assumption. In the regime $N^{1/2}T^4\to0$, where the noise
dominates, one may expect $\hat\theta$ to be asymptotically normal with the
covariance used in Sect.~\ref{estimator}, but we have not proved this, and
no asymptotic statement bears on the coverage at the values of $N$ and $T$
of the tests.

\emph{The constants.} $C_1$ and $T_1$ depend on $\mu$ and on bounds for the
derivatives of the metric on $U$ up to the order that controls the sixth
derivatives of $s$. They are not computed, so Proposition~1 gives no bias
bound at a given window, such as the allowance of 2\% in \eqref{ball}. It
is not uniform over spacetimes without a bound on these derivatives, in
agreement with the obstruction of Sect.~\ref{limits}.

\emph{Further restrictions.} The tips are given events and are not chosen
from the causal set. Vacuum is essential: Ricci terms change the clocks
\eqref{clock} at relative order $T^2$, the order of the signal. The
implementation discards eigenvalues of the clock correlation matrix below
$10^{-8}$ times the largest. Lemma~4 covers this only when no eigenvalue is
discarded; otherwise the lower bound in (A5) has to be replaced by a lower
bound on $J^{\!\top}\widehat WJ$, which is all that the proof uses. The rate
is an upper bound for this estimator under these assumptions, and we claim
no optimality.

\subsection{An idealized random design}\label{app:random}
The conditioning assumption can be removed, in probability, for a
simplified design in which the landmarks are placed at random with a
bounded density. This is not the order-based rule of
Sect.~\ref{estimator}.

\medskip
\noindent\textbf{Corollary 2.} \emph{Let $B_-,B_0,B_+$ be compact subsets
of the open flat interval $\{|x^0|+|\vec x|<\tfrac12\}$ such that $x'-x$ is
future-directed timelike with respect to $\eta$ whenever
$x\in B_-,x'\in B_0$, or $x\in B_0,x'\in B_+$, or $x\in B_-,x'\in B_+$.
Assume (A1) and (A2), and let the landmarks be random, with at least five
in each of the three sets $\exp_o(TB_-)$, $\exp_o(TB_0)$, $\exp_o(TB_+)$,
their numbers being fixed. Assume that their normal coordinates in units
of $T$, with respect to a fixed frame at $o$, have a joint density bounded
by a constant independent of $T$. Let the clocks be all pairs of landmarks
from different sets and all pairs of a landmark and a tip, and let
$\widehat W$ be the weight of Lemma~4 or the identity. Then for every
$\varepsilon>0$ there are $r,C,T_\varepsilon,N_\varepsilon$ such that, for
$T\le T_\varepsilon$ and $N\ge N_\varepsilon$, statements (i) and (ii) of
Proposition~1 hold with $r_1=r$ and $C_1=C$ with probability at least
$1-\varepsilon$. In particular
$\hat\theta-\theta^*_T=O_p(T^4+N^{-1/2})$.}

\noindent\emph{Proof.} Let $\tilde x=(\tilde x_1,\dots,\tilde x_L)$ be the
coordinates in the fixed frame; they range over a compact product
$\mathcal S$ of the three sets. Let $\tilde J(\tilde x)$ be the Jacobian of
the model with respect to all $4L$ coordinates and $(c,d)$, at $c=d=0$. Its
entries are polynomials in $\tilde x$. The model is invariant under a
common spatial rotation of the points and of the tensors, so the three
infinitesimal rotations of the points lie in the kernel of $\tilde J$, and
$\tilde J$ has the same rank at $\tilde x$ and at the rotated configuration
$x$ that satisfies the gauge of (A3). When $x_1^1\neq0$ and $x_2^2\neq0$,
the three gauge conditions define a complement of the infinitesimal
rotations, and therefore $J(\xi,0,0)$ has rank $D$ if and only if
$\tilde J(\tilde x)$ has rank $4L+31$. The configurations of smaller rank
are the common zeros of the minors of order $4L+31$. These polynomials do
not all vanish identically: for five landmarks in each set this is the
exact computation quoted in Sect.~\ref{app:open}, and a further landmark
adds four columns, which are nonzero only in the rows of its own clocks,
where they equal $\pm2\eta(x_{\rm new}-x_j)$ and have rank four for a
suitable position, so that the rank increases by four. Hence the
configurations of smaller rank, together with those in which two points
coincide or the spatial parts of $\tilde x_1$ and $\tilde x_2$ are linearly
dependent, form a set $\mathcal Z$ of Lebesgue measure zero.

For $\mu>0$ let $\mathcal G_\mu\subset\mathcal S$ be the set on which the
singular value of order $4L+31$ of $\tilde J$, the distances between the
points, $|\vec{\tilde x}_1|$ and the area spanned by $\vec{\tilde x}_1$
and $\vec{\tilde x}_2$ are all at least $\mu$. These sets are compact,
they increase to $\mathcal S\setminus\mathcal Z$ as $\mu\to0$, and the
gauge rotation is continuous on them. Their images $\mathcal K_\mu$ in the
gauge of (A3) are compact and satisfy (A3) and (A4), and for small $T$ the
landmarks lie in $I(p,q)$ by Lemma~1. Since the density is bounded, the
probability that $\tilde x\notin\mathcal G_\mu$ is at most a constant times
the Lebesgue measure of $\mathcal S\setminus\mathcal G_\mu$, which tends to
zero as $\mu\to0$, uniformly in $T$. Given $\varepsilon$, choose $\mu$ such
that this probability is at most $\varepsilon/2$, and apply Proposition~1
and Lemma~4 with $\mathcal K=\mathcal K_\mu$, conditionally on the
landmarks. \hfill$\square$

\medskip
The constants of Corollary~2 depend on $\varepsilon$ through $\mu$ and are
not controlled; the remarks of Sect.~\ref{app:open} on the optimizer, the
coverage and the constants apply unchanged.

\subsection{Relation to the theorem of Madsen}
Theorem~4.9 and Corollary~4.10 of Ref.~\cite{madsen} hold for a smooth
globally hyperbolic spacetime of any dimension $d\ge2$, without field
equations. The input is a Poisson sprinkling of known density in a fixed
compact causal diamond. Anchors are selected from the sample by tests on
counts, and the existence of accepted anchors is proved. The output is the
full Riemann tensor at every sample point that passes the tests, in a
reconstructed frame, and the bound
$C[\log(e+\rho)/\rho]^{1/[2(3d+4)]}$ holds simultaneously at all these
points with probability at least $1-\rho^{-n}$ for every $n$, with a
constant that depends on $n$, and almost surely for large $\rho$.

Proposition~1 is weaker in each of these respects. It is restricted to
$d=4$ and to vacuum, the tips are given, the conditioning of the landmarks
and the choice of the local solution are assumed, the tensor is estimated
at one point, and the convergence is in probability. Under these
assumptions the exponent is $1/6$ instead of $1/32$. The two exponents do
not refer to the same problem, and the comparison does not show that one
construction is more efficient than the other. In particular, the rate
$\rho^{-1/6}$ relies on vacuum, which makes the volume clocks accurate to
relative order $T^4$ (Lemma~2); without it the clocks are accurate only to
relative order $T^2$ and the present estimator is not consistent.

\subsection{Numerical check of the orders}
The script \texttt{check\_convergence\_orders.py} in the repository checks
the orders used above on exact rotating plane waves
(Sect.~\ref{tests}), for a fixed configuration of 48 landmarks in windows
scaled by $1$, $\tfrac12$, $\tfrac14$ and $\tfrac18$ and with exact
interval volumes in place of counts. Over the three successive halvings of
the window, the residual of the fifth-degree expansion of Lemma~1(a)
decreases by factors of $15.1$, $15.8$ and $15.9$, and the relative volume
correction of Lemma~2 by $15.8$, $15.9$ and $16.0$, to be compared with
$16$. The relative error of the fitted tensor decreases by $4.3$, $4.1$
and $4.0$ without $\sigma_2$ and by $3.8$, $4.0$ and $4.0$ with it, to be
compared with $4$, and the error $|\hat\theta-\theta^*_T|$ of all
parameters by $22.8$, $18.0$ and $16.6$ without $\sigma_2$ and by $26.9$,
$17.4$ and $16.6$ with it, to be compared with $16$. The same script
verifies the statement on Myrheim's formula made after Lemma~2 and
performs the exact rank computation
quoted in Sect.~\ref{app:open}. This checks the orders in one family of
spacetimes and one configuration; it does not test the assumptions.
 in note.tex.
% Packages: amsmath and amssymb only (already loaded by note.tex); no amsthm.
% Labels of note.tex used: setting, estimator, tests, limits, clock, cov, model, square, ball.
% Citation keys used: brewin, gibbons, wang, madsen, and ONE new key, oneill:
%   \bibitem{oneill} O'Neill, B.: Semi-Riemannian Geometry with Applications to
%   Relativity. Academic Press, New York (1983)
% Check of the orders and of the rank certificate: python check_convergence_orders.py
\section{A conditional convergence statement}\label{app:conv}
This appendix states what we can prove about the fit of
Sect.~\ref{estimator} in shrinking windows, and what we cannot. The
statement is conditional: it assumes that the landmarks form a
well-conditioned configuration, and it concerns the solution of the fit that
lies near the true geometry. It does not show that the selection rule and
the optimizer of Sect.~\ref{estimator} meet these two conditions, and it
says nothing about the coverage of the ball \eqref{ball}. Only orders are
established; no constant is evaluated.

\subsection{Setting and assumptions}\label{app:setting}
We write $|\cdot|$ for the Euclidean norm of a vector of components,
$\|\cdot\|$ for the corresponding operator norm, and $A\succeq B$ if $A-B$
is positive semidefinite. For two events $x,x'$ of a convex normal
neighbourhood, $s(x,x')=g(v,v)$ with $v=\exp_x^{-1}x'$ is the signed squared
geodesic interval, so that $s=-\tau^2$ for timelike separation, and the
quantity $\sigma(x,y)$ of Sect.~\ref{estimator} is $s/T^2$.

\begin{description}
\item[(A1) Geometry.] $(M,g)$ is a time-oriented four-dimensional Lorentzian
manifold with smooth metric, $o\in M$, and $(e_a)$ is an orthonormal frame
at $o$ with $e_0$ future-directed. $U$ is a convex normal neighbourhood of
$o$, that is, an open set that is a normal neighbourhood of each of its
points; then any two points of $U$ are joined by a unique geodesic segment
in $U$, and $(x,x')\mapsto\exp_x^{-1}x'$ is smooth on $U\times U$
\cite[Chap.~5]{oneill}. The metric is Ricci-flat on $U$, and in the normal
coordinates based at $(o,e_a)$ its matrix of components satisfies
$\|g_{ab}(x)-\eta_{ab}\|\le1/10$ on $U$. Every point of a Ricci-flat region
has arbitrarily small neighbourhoods of this kind. For $0<T\le T_0$ the tips
are $p=\exp_o(-\tfrac T2e_0)$ and $q=\exp_o(\tfrac T2e_0)$, and, as in
Sect.~\ref{setting}, $I(p,q)=J^+(p)\cap J^-(q)\subset U$; this holds for
small $T$ if $M$ is strongly causal at $o$. The proofs use finitely many
derivatives, six of $s$ and four of the metric components in normal
coordinates; we have not determined the least differentiability of $g$ that
provides them.
\item[(A2) Sampling.] Conditionally on the landmarks, the $N$ measurement
events are independent and uniformly distributed on $I(p,q)$ with respect to
the volume of $g$. This is the case when the landmarks are taken from an
independent selection sample.
\item[(A3) Landmark configuration.] There are $L$ landmarks
$\ell_1,\dots,\ell_L\in I(p,q)$; put $\ell_0=p$ and $\ell_{L+1}=q$. Their
normal coordinates in units of $T$ are $x_i\in\mathbb R^4$, that is,
$\ell_i=\exp_o(Tx_i^ae_a)$, so that $x_0=(-\tfrac12,0,0,0)$ and
$x_{L+1}=(\tfrac12,0,0,0)$. The spatial axes are those of the gauge of
Sect.~\ref{estimator}: $x_1=(x_1^0,x_1^1,0,0)$ with $x_1^1>0$ and
$x_2=(x_2^0,x_2^1,x_2^2,0)$ with $x_2^2>0$. This fixes the frame up to the
reflection of $e_3$, and we fix one of the two frames. The remaining $4L-3$
coordinates form $\xi\in\mathbb R^{4L-3}$. The clocks are the pairs
$(\ell_i,\ell_j)$ with $(i,j)$ in a fixed set $\mathcal P$ of $m$ index
pairs, $(i,j)\neq(0,L+1)$. We assume that $\xi$ lies in a fixed compact set
$\mathcal K\subset\mathbb R^{4L-3}$ such that, for every configuration in
$\mathcal K$, the $L+2$ points $x_i$ are distinct, $x_1,\dots,x_L$ lie in
the open flat interval $\{|x^0|+|\vec x|<\tfrac12\}$, and $x_j-x_i$ is
future-directed timelike with respect to $\eta$ for every
$(i,j)\in\mathcal P$.
\item[(A4) Conditioning.] Let $\theta=(\xi,c,d)\in\mathbb R^D$, $D=4L+31$,
where $c\in\mathbb R^{10}$ and $d\in\mathbb R^{24}$ are the components of the
curvature and of its gradient in the bases of Sect.~\ref{estimator} (there
$d$ is called $g$, a letter needed here for the metric). Let
$\sigma(\theta)\in\mathbb R^m$ be the model \eqref{model} for the pairs of
$\mathcal P$, with $\sigma_2$ as in \eqref{square} or with $\sigma_2=0$, and
$J(\theta)=\partial\sigma/\partial\theta$. We assume that the $m\times D$
matrix $J(\xi,0,0)$ has rank $D$ for every $\xi\in\mathcal K$.
\item[(A5) Weights.] The fit minimizes
$\chi^2(\theta)=N[y-\sigma(\theta)]^{\!\top}\widehat W[y-\sigma(\theta)]$
with a symmetric matrix $\widehat W$ that may depend on the data, and there
are constants $0<w_-\le w_+$ such that $w_-I\preceq\widehat W\preceq w_+I$
with probability at least $1-\varpi_N$, where $\varpi_N\to0$ as
$N\to\infty$. For the weight of Sect.~\ref{estimator},
$\widehat W=(N\widehat\Sigma)^{-1}$, this is proved in
Lemma~4 below.
\end{description}

Assumptions (A3) and (A4) are substantive, and (A4) is a condition on the
shape of the configuration only. By compactness there are $\tau_0>0$ and
$\mu>0$ such that $\eta(x_j-x_i,x_j-x_i)\le-\tau_0^2$ for
$(i,j)\in\mathcal P$ and the smallest singular value of $J(\xi,0,0)$ is at
least $\mu$, for all $\xi\in\mathcal K$; all constants below depend on
these numbers. The spatial axes of the frame $(e_a)$ depend on the
landmarks; since the orthonormal frames at $o$ with the given $e_0$ form a
compact set, all estimates below are uniform over them. The true parameter
is $\theta^*_T=(\xi,c^*,d^*)$, where $\xi$
is the true configuration and $c^*$, $d^*$ are the components of
$T^2C_{abcd}(o)$ and of $T^3\nabla_eC_{abcd}(o)$ in the frame $(e_a)$. The
measured clocks are $y_{ij}=-(n_{ij}/N)^{1/2}$, with $n_{ij}$ the number of
measurement events in $I(\ell_i,\ell_j)$, and the continuum clocks are
$\bar y_{ij}=-[V(\ell_i,\ell_j)/V(p,q)]^{1/2}$.

\subsection{Statement}
\noindent\textbf{Proposition 1.} \emph{Assume (A1)--(A5). There are
constants $T_1,r_1>0$ and $C_1,C_2<\infty$, which depend only on the
geometry of $(U,g)$, on $\mathcal K$, $\mathcal P$ and on $w_\pm$, such that
for all $0<T\le T_1$, all $N\ge2$ and every landmark configuration allowed
by (A3) the following holds with probability at least $1-\varpi_N-C_2/N$.}
\begin{enumerate}
\item[(i)] \emph{$\chi^2$ has exactly one critical point
$\hat\theta=(\hat\xi,\hat c,\hat d)$ in the closed ball of radius $r_1$
around $\theta^*_T$, and $\hat\theta$ is the unique minimizer of $\chi^2$
on this ball.}
\item[(ii)] \emph{$|\hat\theta-\theta^*_T|\le C_1\left(T^4+|y-\bar y|\right)$.}
\end{enumerate}
\emph{Moreover $\operatorname{E}|y-\bar y|^2\le m/N$. Consequently, for
every $t>0$, with probability at least $1-\varpi_N-C_2/N-m/t^2$,}
\begin{equation}
 |\hat\theta-\theta^*_T|\le C_1\left(T^4+tN^{-1/2}\right),\qquad
 \left|\hat c/T^2-c^*/T^2\right|\le C_1\left(T^2+tN^{-1/2}T^{-2}\right),
 \label{conv:rate}
\end{equation}
\emph{where $c^*/T^2$ are the components of the Weyl tensor $C_{abcd}(o)$.
The same holds, with other constants, in the tidal norm of
Sect.~\ref{setting}.}

\medskip
In short, $\hat\theta-\theta^*_T=O(T^4)+O_p(N^{-1/2})$ in the normalized
variables of the fit. The statement and its proof are the same with and
without $\sigma_2$, because $\sigma_2=O(|c|^2)$; in particular the
coefficient of the volume term in \eqref{square} is not used. The estimator
returns $\hat c$, an estimate of $T^2C(o)$. The second inequality of
\eqref{conv:rate} refers to the tensor in physical units, which requires
$T$, or equivalently bounds the relative error when $C(o)\neq0$.

\medskip
\noindent\textbf{Corollary 1.} \emph{Let the measurement events be a
Poisson sprinkling of density $\rho$ in $I(p,q)$, with lengths measured in
a fixed unit, and let (A1)--(A5) hold. If $T\to0$ and $\rho T^8\to\infty$,
then the solution $\hat\theta$ of Proposition~1 exists with probability
tending to one, and}
\begin{equation}
 \hat c/T^2-C(o)=O(T^2)+O_p\!\left(\rho^{-1/2}T^{-4}\right)\to0
 \quad\text{in probability}.
\end{equation}
\emph{For $T$ proportional to $\rho^{-1/12}$ the error is
$O_p(\rho^{-1/6})$.}

\subsection{Proofs}
\noindent\textbf{Lemma 1} (normal coordinates). \emph{Let $\mathcal F$ be a
compact set of orthonormal frames $(x,f)$ with $x\in U$. Let
$w\in\mathbb R^4$ be the normal coordinates based at $(x,f)$, and let
$s_{x,f}(w,w')$ and $\gamma_{x,f}(w)$ be the function $s$ and the matrix of
metric components in these coordinates. There are $r_0>0$ and $C_0$ such
that for all $(x,f)\in\mathcal F$ and $|w|,|w'|\le r_0$, with
$\Delta=w-w'$:}
\begin{enumerate}
\item[(a)] \emph{$s_{x,f}(w,w')=\eta(\Delta,\Delta)-\tfrac13R(w,w',w,w')
-\tfrac1{12}(w+w')^e\nabla_eR(w,w',w,w')+\varrho$ with
$|\varrho|\le C_0(|w|+|w'|)^6$, where $R$ and $\nabla R$ are the components
of the curvature and its gradient at $x$ in the frame $f$;}
\item[(b)] \emph{$s_{x,f}(w,w')-\eta(\Delta,\Delta)$ is bounded by
$C_0(|w|+|w'|)^4$ and its first derivatives by $C_0(|w|+|w'|)^3$;}
\item[(c)] $\bigl|\,|\det\gamma_{x,f}(w)|^{1/2}-1\,\bigr|\le C_0|w|^4$;
\item[(d)] $s_{x,f}(0,w)=\eta(w,w)$.
\end{enumerate}

\noindent\emph{Proof.} The maps $(x,f,w,w')\mapsto s_{x,f}(w,w')$ and
$(x,f,w)\mapsto\gamma_{x,f}(w)$ are smooth on an open set that contains
$\mathcal F\times\{0\}\times\{0\}$, hence on
$\mathcal F\times\{|w|,|w'|\le r_0\}$ for some $r_0>0$. Statement (d) is the
definition of $s$. For (a) apply Taylor's theorem in $(w,w')$ at the
origin: the derivatives of order six are bounded on this compact set, and
the Taylor polynomial of degree five is the expansion of the squared
geodesic interval in normal coordinates, Sect.~9 of Ref.~\cite{brewin}
(Eqs.~(11.21a)--(11.21c) in arXiv:0903.2087v3). That polynomial has no terms
of degree three, which gives (b) in the same way. For (c),
$\gamma_{ab}(w)=\eta_{ab}-\tfrac13R_{acbd}w^cw^d
-\tfrac16\nabla_eR_{acbd}w^cw^dw^e+O(|w|^4)$ \cite{brewin}, and the traces
of the second and third terms with $\eta^{ab}$ contain only the Ricci tensor
and its gradient, which vanish; hence $\det\gamma=-1+O(|w|^4)$, with a
remainder bounded as before. \hfill$\square$

\medskip
\noindent\textbf{Lemma 2} (volume of a small vacuum interval). \emph{Let
$\mathcal T$ be a compact set of future-directed unit timelike vectors
$(x,u)$ with $x\in U$. There are $\tau_1>0$ and $C_V$ such that, whenever
$(x,u)\in\mathcal T$, $0<\tau\le\tau_1$, $x'=\exp_x(\tau u)$ and
$J^+(x)\cap J^-(x')\subset U$,}
\begin{equation}
 \left|V(x,x')-\tfrac{\pi}{24}\tau^4\right|\le C_V\tau^8 .
 \label{conv:volume}
\end{equation}

\noindent\emph{Proof.} Complete $u$ to an orthonormal frame $f$ with
$f_0=u$. The frames so obtained form a compact set, to which Lemma~1
applies. We use the normal coordinates $w$ based at $(x,f)$, in which $x'$
is $\tau\hat e$ with $\hat e=(1,0,0,0)$, write $s=s_{x,f}$, and put
$\Xi(w)=s(0,w)-s(w,\tau\hat e)$. We show that the half
$H_-=\{z\in J^+(x)\cap J^-(x'):\Xi\ge0\}$ of the interval has volume
$\tfrac{\pi}{48}\tau^4+O(\tau^8)$. The other half, where $\Xi\le0$, is
treated in the same way from $x'$ with the time orientation reversed, and
the common boundary $\Xi=0$ has measure zero. All $O(\cdot)$ below are
uniform over $\mathcal T$ and over the directions $\nu$.

(1) \emph{Localization.} By (A1), every causal vector $v$ at a point of $U$
has $|\vec v|\le\tfrac54|v^0|$ in the coordinates of (A1), with $v^0>0$ if
it is future-directed. Hence $x^0$ increases along future-directed causal
curves in $U$, and the coordinate differences satisfy
$|z-x|\le\tfrac94|x'-x|$ for every $z\in J^+(x)\cap J^-(x')$. Since the
change to the coordinates $w$ is smooth, uniformly over the frames, there
is $C\ge1$ such that $|w|\le C\tau$ for all these $z$.

(2) \emph{Rays.} In a convex normal neighbourhood, $z\in J^+(x)$ if and
only if $\exp_x^{-1}z$ is zero or future-directed causal
\cite[Lemma~14.2]{oneill}, that is, if and only if $w=t\nu$ with $t\ge0$
and $\nu\in\mathcal N=\{(1,\vec\nu):|\vec\nu|\le1\}$. For $\nu\in\mathcal N$
let $\phi_\nu(t)=\Xi(t\nu)$. By Lemma~1(d),
$\phi_\nu(t)=t^2\eta(\nu,\nu)-s(t\nu,\tau\hat e)$, and
$\eta(t\nu-\tau\hat e,t\nu-\tau\hat e)=t^2\eta(\nu,\nu)+2\tau t-\tau^2$.
By Lemma~1(b), for $0\le t\le C\tau$,
\begin{equation}
 \phi_\nu(t)=\tau^2-2\tau t+O(\tau^4),\qquad
 \phi_\nu'(t)=-2\tau+O(\tau^3).
\end{equation}
For $\tau\le\tau_1$ small, $\phi_\nu$ is therefore strictly decreasing on
$[0,C\tau]$, with $\phi_\nu(0)=\tau^2$ and $\phi_\nu(\tau)<0$. It has a
single zero $t_*(\nu)\in(0,\tau)$, which depends continuously on $\nu$.

(3) \emph{The half-interval is
$\{t\nu:\nu\in\mathcal N,\ 0\le t\le t_*(\nu)\}$.} If $z\in H_-$, then
$w=t\nu$ with $t\le C\tau$ by (1) and (2), and $\phi_\nu(t)\ge0$ gives
$t\le t_*(\nu)$. Conversely, let $0\le t\le t_*(\nu)$. On the segment from
$0$ to $t\nu$ we have $\Xi\ge0$, that is, $s(\cdot,x')\le s(x,\cdot)\le0$,
so the segment lies in the set of events causally related to $x'$, and it
does not contain $x'$, where $\Xi=-\tau^2$. The sets
$J^-(x')\setminus\{x'\}$ and $J^+(x')\setminus\{x'\}$ are disjoint and
closed in $U\setminus\{x'\}$, and the segment starts at $x\in J^-(x')$;
being connected, it lies in $J^-(x')$. Hence $t\nu\in H_-$.

(4) \emph{Volume.} With $w=(t,t\vec\nu)$ we have
$\mathrm d^4w=t^3\,\mathrm dt\,\mathrm d^3\vec\nu$, and Lemma~1(c) gives
\begin{equation}
 \operatorname{Vol}(H_-)=\int_{|\vec\nu|\le1}\int_0^{t_*(\nu)}
 \left[1+O(\tau^4)\right]t^3\,\mathrm dt\,\mathrm d^3\vec\nu
 =\frac14\int_{|\vec\nu|\le1}t_*(\nu)^4\,\mathrm d^3\vec\nu+O(\tau^8).
 \label{conv:half}
\end{equation}

(5) \emph{Expansion of $t_*$.} Insert Lemma~1(a) with $w=t\nu$ and
$w'=\tau\hat e$. Since
$R(t\nu,\tau\hat e,t\nu,\tau\hat e)=t^2\tau^2E_{ij}\nu^i\nu^j$ with
$E_{ij}=R_{0i0j}$,
\begin{equation}
 \phi_\nu(t)=\tau^2-2\tau t+\tfrac13t^2\tau^2E_{ij}\nu^i\nu^j
 +\tfrac1{12}t^2\tau^2\left[(t+\tau)\nabla_0E_{ij}
 +t\,\nu^k\nabla_kE_{ij}\right]\nu^i\nu^j+O(\tau^6).
\end{equation}
The zero satisfies $t_*=\tfrac\tau2+O(\tau^3)$, and then
\begin{equation}
 t_*(\nu)=\frac\tau2\left[1+\frac{\tau^2}{12}E_{ij}\nu^i\nu^j
 +\frac{\tau^3}{96}\left(3\nabla_0E_{ij}+\nu^k\nabla_kE_{ij}\right)
 \nu^i\nu^j+O(\tau^4)\right].
\end{equation}
The integral of $\nu^i\nu^j$ over the unit ball is
$\tfrac{4\pi}{15}\delta^{ij}$, and that of $\nu^i\nu^j\nu^k$ vanishes. In
vacuum $E_{ii}=R_{00}=0$ and $\nabla_0E_{ii}=\nabla_0R_{00}=0$, so the terms
of relative order $\tau^2$ and $\tau^3$ integrate to zero, and
\eqref{conv:half} gives
$\operatorname{Vol}(H_-)=\tfrac14\cdot\tfrac{4\pi}3(\tau/2)^4+O(\tau^8)
=\tfrac{\pi}{48}\tau^4+O(\tau^8)$. \hfill$\square$

\medskip
Lemma~2 uses the field equations only through $R_{ab}=0$, and it gives only
the order of the correction. As a check of the method, without that
condition the same computation, with the volume element
$1-\tfrac16R_{ab}w^aw^b$, gives
$V=\tfrac{\pi}{24}\tau^4[1-\tfrac1{180}R^a{}_a\tau^2+\tfrac1{30}R_{00}\tau^2
+O(\tau^3)]$, which is the formula of Myrheim, Eqs.~(7), (14) and (23) of
Ref.~\cite{gibbons}. In vacuum the first correction is of relative order
$\tau^4$, in agreement with Eq.~(79) of Ref.~\cite{wang}; its coefficient, on
which \eqref{square} and Ref.~\cite{wang} differ, plays no role here.

\medskip
\noindent\textbf{Lemma 3} (clocks). \emph{Assume (A1)--(A3). (a) There are
$T_1>0$ and $C_b$ such that $|\bar y_{ij}-\sigma_{ij}(\theta^*_T)|\le
C_bT^4$ for all $T\le T_1$, $\xi\in\mathcal K$ and $(i,j)\in\mathcal P$.
(b) $\operatorname{E}|y-\bar y|^2\le m/N$.}

\noindent\emph{Proof.} (a) Lemma~1(a) at $(o,e)$ with $w=Tx_i$ and
$w'=Tx_j$, and $R_{abcd}=C_{abcd}$ in vacuum, give
\begin{equation}
 \frac{s(\ell_i,\ell_j)}{T^2}=\eta(\Delta,\Delta)
 -\frac{T^2}3C(x_i,x_j,x_i,x_j)
 -\frac{T^3}{12}(x_i+x_j)^e\nabla_eC(x_i,x_j,x_i,x_j)+O(T^4),
\end{equation}
with $\Delta=x_i-x_j$ and the tensors at $o$. The right-hand side is the
model \eqref{model} at $\theta^*_T$ without $\sigma_2$, up to $O(T^4)$, and
$\sigma_2(\theta^*_T)=O(|c^*|^2)=O(T^4)$. Since
$\eta(\Delta,\Delta)\le-\tau_0^2$, for small $T$ we have
$\tau_{ij}^2=-s(\ell_i,\ell_j)\ge\tau_0^2T^2/2$ and $\tau_{ij}=O(T)$. In
the coordinates of (A1) the vector $\exp_{\ell_i}^{-1}\ell_j$ has components
$T(x_j-x_i)+O(T^2)$, so for small $T$ it is future-directed and $\ell_j$
lies in the timelike future of $\ell_i$. The unit tangent at $\ell_i$ of the
geodesic to $\ell_j$ then has bounded components, and these vectors lie in a
compact set $\mathcal T$.
Moreover $J^+(\ell_i)\cap J^-(\ell_j)\subset I(p,q)\subset U$. Lemma~2
therefore gives $V(\ell_i,\ell_j)=\tfrac{\pi}{24}\tau_{ij}^4[1+O(T^4)]$, and
likewise $V(p,q)=\tfrac{\pi}{24}T^4[1+O(T^4)]$. Hence
$\bar y_{ij}=-(\tau_{ij}^2/T^2)[1+O(T^4)]=s(\ell_i,\ell_j)/T^2+O(T^4)$.

(b) Given the landmarks, $n_{ij}$ is binomial with parameters $N$ and
$\pi_{ij}=V(\ell_i,\ell_j)/V(p,q)=\bar y_{ij}^2$. Since
$|\sqrt a-\sqrt b|\le|a-b|/\sqrt b$,
\begin{equation}
 |y_{ij}-\bar y_{ij}|\le\frac{|n_{ij}/N-\pi_{ij}|}{\sqrt{\pi_{ij}}},
 \qquad\operatorname{E}|y_{ij}-\bar y_{ij}|^2\le\frac{1-\pi_{ij}}N .
\end{equation}
Summing over the $m$ pairs gives (b). The clocks are strongly correlated,
see \eqref{cov}, but no independence is used. \hfill$\square$

\medskip
\noindent\textbf{Lemma 4} (weights). \emph{Let $N\widehat\Sigma=\widehat
D\,\widehat\Gamma\,\widehat D$ be the covariance of Sect.~\ref{estimator},
where $\widehat\Gamma$ is the empirical covariance matrix of the membership
indicators of the $m$ intervals and
$\widehat D=\operatorname{diag}(\tfrac12\hat\pi_{ij}^{-1/2})$ with
$\hat\pi_{ij}=n_{ij}/N$. Under (A1)--(A3) there are $T_2>0$,
$0<w_-\le w_+$ and $C_W$ such that $\widehat W=(N\widehat\Sigma)^{-1}$
satisfies (A5) with $\varpi_N=C_W/N$ for all $T\le T_2$.}

\noindent\emph{Proof.} Let $\Gamma_T(\xi)$ be the covariance matrix of the
indicators under the sampling law, and $\Gamma_0(\xi)$ the same matrix for
the flat metric $\eta$ and the points $x_i$.

(1) \emph{$\Gamma_0(\xi)$ is positive definite.} Otherwise a combination
$\sum\lambda_{ij}1_{ij}$ of the indicators with $\lambda\neq0$ is constant
almost everywhere in the flat interval. Fix a pair $(i,j)$. The edge of its
interval, where the future cone of $x_i$ meets the past cone of $x_j$, is a
round two-sphere, and the only null cones that contain it are those of
$x_i$ and $x_j$; the null cone of any other point meets it in the zero set
of a quadratic polynomial that does not vanish identically on the sphere,
a closed set without interior points. Choose $z$ on the edge and on none of
the other cones, and a ball $B$ around $z$ that meets no null cone of a
point $x_k$ other than the future cone of $x_i$ and the past cone of $x_j$.
In $B$, every indicator other than $1_{ij}$ is a constant, or a constant
times the indicator of $J^+(x_i)$ or of $J^-(x_j)$, while $1_{ij}$ is the
product of these two. The two cones cross transversally at $z$. If
$1\le i,j\le L$, all four combinations of the two indicators occur on sets
of positive measure in $B$, and constancy of the sum forces
$\lambda_{ij}=0$. If $i=0$, only $B\cap J^+(p)$ carries probability, and no
interval $I(x_k,x_j)$ with $k\neq0$ meets $B$, because
$z\in\partial J^+(p)$ whereas $x_k$ lies in the open interval; the two
values of the indicator of $J^-(x_j)$ again force $\lambda_{0j}=0$. The
case $j=L+1$ is symmetric.

(2) In the coordinates $x/T$ the metric has components $g_{ab}(Tx)$. These,
the inverse exponential map and the volume density converge to their flat
values as $T\to0$, uniformly on compact sets. The indicator of each
interval is therefore continuous in $(T,\xi)$ and in the spatial frame at
$o$ at every point off the boundary of the limiting interval, and by
dominated convergence $\Gamma_T(\xi)$ is continuous in these variables on
$[0,T_0]\times\mathcal K$ times the compact set of frames, with the value
$\Gamma_0(\xi)$ at $T=0$. With (1) and compactness,
$\Gamma_T(\xi)\succeq4\gamma I$ for some $\gamma>0$, all $\xi\in\mathcal K$,
all frames and all $T\le T_2$. In the same way $\pi_{ij}\ge\pi_->0$.

(3) $\widehat\Gamma$ and $\hat\pi$ are built from averages of bounded
variables, so $\operatorname{E}\|\widehat\Gamma-\Gamma_T\|^2\le Cm^2/N$ and
$\operatorname{E}|\hat\pi-\pi|^2\le m/N$. By Chebyshev's inequality the
event that $\|\widehat\Gamma-\Gamma_T\|\le2\gamma$ and
$\hat\pi_{ij}\ge\pi_-/2$ for all pairs has probability at least $1-C_W/N$.
On this event the diagonal entries of $\widehat D$ lie between $\tfrac12$
and $(2\pi_-)^{-1/2}$, $\widehat\Gamma\succeq2\gamma I$ and
$\operatorname{tr}\widehat\Gamma\le m/2$, hence
$\tfrac\gamma2I\preceq N\widehat\Sigma\preceq\tfrac{m}{4\pi_-}I$.
\hfill$\square$

\medskip
\noindent\textbf{Lemma 5} (local minimizer). \emph{Let $F$ be a twice
continuously differentiable map from the closed ball of radius $r_0$ around
$\theta^*\in\mathbb R^D$ to $\mathbb R^m$, with $\|DF\|\le M_1$ and
$\|D^2F\|\le M_2$. Let $W$ be symmetric with $0\preceq W\preceq w_+I$ and
$DF(\theta^*)^{\!\top}W\,DF(\theta^*)\succeq\lambda I$, $\lambda>0$. For
$y\in\mathbb R^m$ let $Q(\theta)=\tfrac12[y-F(\theta)]^{\!\top}
W[y-F(\theta)]$ and $\epsilon=|y-F(\theta^*)|$, and put}
\begin{equation}
 r_1=\min\left\{r_0,\frac{\lambda}{12M_1M_2w_+}\right\},\qquad
 \epsilon_1=\min\left\{\frac{\lambda}{4M_2w_+},
 \frac{\lambda r_1}{8M_1w_+}\right\}.
\end{equation}
\emph{If $\epsilon\le\epsilon_1$, then $Q$ has exactly one critical point
$\hat\theta$ in the closed ball of radius $r_1$ around $\theta^*$, it is the
unique minimizer of $Q$ on this ball, and
$|\hat\theta-\theta^*|\le4M_1w_+\epsilon/\lambda$.}

\noindent\emph{Proof.} Let $|\theta-\theta^*|\le r_1$, let $v$ be a unit
vector and $\rho=y-F(\theta)$. Then
\begin{equation}
 v^{\!\top}D^2Q(\theta)\,v=[DF(\theta)v]^{\!\top}W[DF(\theta)v]
 -\rho^{\!\top}W\,D^2F(\theta)[v,v].
\end{equation}
Here $|\rho|\le\epsilon+M_1r_1$ and $\|DF(\theta)-DF(\theta^*)\|\le M_2r_1$,
so that
$v^{\!\top}D^2Q(\theta)v\ge\lambda-2M_1M_2w_+r_1-M_2w_+(\epsilon+M_1r_1)
\ge\lambda/2$. Thus $Q$ is strongly convex on the ball. It has a unique
minimizer $\hat\theta$ there and at most one critical point, and
$Q(\theta)\ge Q(\theta^*)+DQ(\theta^*)(\theta-\theta^*)
+\tfrac\lambda4|\theta-\theta^*|^2$. Since
$|DQ(\theta^*)|=|DF(\theta^*)^{\!\top}W[y-F(\theta^*)]|\le M_1w_+\epsilon$,
we get $Q(\theta)>Q(\theta^*)$ whenever
$|\theta-\theta^*|>4M_1w_+\epsilon/\lambda$. Hence
$|\hat\theta-\theta^*|\le4M_1w_+\epsilon/\lambda\le r_1/2$, so $\hat\theta$
is an interior point and therefore a critical point. \hfill$\square$

\medskip
\noindent\emph{Proof of Proposition 1.} The map $\sigma$ is a polynomial in
$\theta$. Let $M_1$ and $M_2$ bound its first and second derivatives on the
set of all $\theta$ at distance at most one from
$\{(\xi,c,d):\xi\in\mathcal K,\ |c|,|d|\le1\}$; for small $T$ this set
contains the ball of radius $r_0=1$ around $\theta^*_T$. Because $\sigma_2$
is quadratic in $c$, the matrix $J(\xi,0,0)$ does not depend on it, and
$\|J(\theta^*_T)-J(\xi,0,0)\|\le M_2(|c^*|+|d^*|)=O(T^2)$. By (A4) the
smallest singular value of $J(\theta^*_T)$ is at least $\mu/2$ for
$T\le T_1$, so that on the event of (A5)
$J(\theta^*_T)^{\!\top}\widehat WJ(\theta^*_T)\succeq\lambda I$ with
$\lambda=w_-\mu^2/4$. Apply Lemma~5 on this event with $F=\sigma$,
$W=\widehat W$ and $\theta^*=\theta^*_T$; the functions $\chi^2$ and $Q$
differ by the factor $2N$. By Lemma~3(a),
$\epsilon=|y-\sigma(\theta^*_T)|\le|y-\bar y|+\sqrt m\,C_bT^4$. Reduce $T_1$
so that $\sqrt m\,C_bT_1^4\le\epsilon_1/2$. By Lemma~3(b) and Chebyshev's
inequality, $|y-\bar y|\le\epsilon_1/2$ except on an event of probability at
most $C_2/N$ with $C_2=4m/\epsilon_1^2$. Lemma~5 then gives (i), and (ii)
with $C_1=4M_1w_+\max(1,\sqrt m\,C_b)/\lambda$. Finally
$|y-\bar y|\le tN^{-1/2}$ except with probability $m/t^2$, which gives
\eqref{conv:rate}, the second inequality because
$|\hat c-c^*|\le|\hat\theta-\theta^*_T|$. \hfill$\square$

\medskip
\noindent\emph{Proof of Corollary 1.} Given their number $N$, the events of
a Poisson sprinkling in $I(p,q)$ satisfy (A2), and $N$ is Poisson with mean
$\rho V(p,q)=\tfrac{\pi}{24}\rho T^4[1+O(T^4)]$ by Lemma~2. If
$\rho T^4\to\infty$, then $N\ge\rho T^4/20$ with probability tending to
one, and on this event $N^{-1/2}T^{-2}\le\sqrt{20}\,\rho^{-1/2}T^{-4}$ and
$\varpi_N+C_2/N\to0$. Insert this into \eqref{conv:rate}. The two terms
balance for $T^6$ proportional to $\rho^{-1/2}$, and then
$\rho T^8\to\infty$. \hfill$\square$

\subsection{What is not proved}\label{app:open}
\emph{The selection rule.} Assumptions (A3) and (A4) are not shown to hold
for the landmarks produced by the rule of Sect.~\ref{estimator}, which are
events of a selection sample chosen by counts, with a clock set chosen by
counts. Assumption (A4) is not vacuous: for an explicit rational
configuration with five landmarks in each of three time bands, and with
the clocks between different bands and to the tips, exact modular
arithmetic gives $J(\xi,0,0)$ the full rank $D=91$. Since full rank is the
non-vanishing of a polynomial, (A4) then holds for almost every
configuration with these clocks (see Corollary~2). This identifies the
parameters; it does not bound $\mu$, on which the constants depend. In
that configuration the singular values of $J(\xi,0,0)$ range over eight
orders of magnitude.

\emph{The optimizer.} Proposition~1 concerns the critical point near
$\theta^*_T$. It does not show that the Levenberg--Marquardt iteration from
the four starts of Sect.~\ref{estimator} converges to this point, and it
does not exclude other local minima. The reflections of
Sect.~\ref{setting} map $\hat\theta$ to further minimizers with the same
$\chi^2$; these are the mirror images accounted for there. No global
identifiability is proved.

\emph{The coverage.} Nothing is proved about the ball \eqref{ball}. For
$T$ proportional to $\rho^{-1/12}$ the bias and the noise in
\eqref{conv:rate} are of the same order, so a ball built from the noise
covariance alone cannot be expected to cover, and the allowance $\delta$
remains an assumption. In the regime $N^{1/2}T^4\to0$, where the noise
dominates, one may expect $\hat\theta$ to be asymptotically normal with the
covariance used in Sect.~\ref{estimator}, but we have not proved this, and
no asymptotic statement bears on the coverage at the values of $N$ and $T$
of the tests.

\emph{The constants.} $C_1$ and $T_1$ depend on $\mu$ and on bounds for the
derivatives of the metric on $U$ up to the order that controls the sixth
derivatives of $s$. They are not computed, so Proposition~1 gives no bias
bound at a given window, such as the allowance of 2\% in \eqref{ball}. It
is not uniform over spacetimes without a bound on these derivatives, in
agreement with the obstruction of Sect.~\ref{limits}.

\emph{Further restrictions.} The tips are given events and are not chosen
from the causal set. Vacuum is essential: Ricci terms change the clocks
\eqref{clock} at relative order $T^2$, the order of the signal. The
implementation discards eigenvalues of the clock correlation matrix below
$10^{-8}$ times the largest. Lemma~4 covers this only when no eigenvalue is
discarded; otherwise the lower bound in (A5) has to be replaced by a lower
bound on $J^{\!\top}\widehat WJ$, which is all that the proof uses. The rate
is an upper bound for this estimator under these assumptions, and we claim
no optimality.

\subsection{An idealized random design}\label{app:random}
The conditioning assumption can be removed, in probability, for a
simplified design in which the landmarks are placed at random with a
bounded density. This is not the order-based rule of
Sect.~\ref{estimator}.

\medskip
\noindent\textbf{Corollary 2.} \emph{Let $B_-,B_0,B_+$ be compact subsets
of the open flat interval $\{|x^0|+|\vec x|<\tfrac12\}$ such that $x'-x$ is
future-directed timelike with respect to $\eta$ whenever
$x\in B_-,x'\in B_0$, or $x\in B_0,x'\in B_+$, or $x\in B_-,x'\in B_+$.
Assume (A1) and (A2), and let the landmarks be random, with at least five
in each of the three sets $\exp_o(TB_-)$, $\exp_o(TB_0)$, $\exp_o(TB_+)$,
their numbers being fixed. Assume that their normal coordinates in units
of $T$, with respect to a fixed frame at $o$, have a joint density bounded
by a constant independent of $T$. Let the clocks be all pairs of landmarks
from different sets and all pairs of a landmark and a tip, and let
$\widehat W$ be the weight of Lemma~4 or the identity. Then for every
$\varepsilon>0$ there are $r,C,T_\varepsilon,N_\varepsilon$ such that, for
$T\le T_\varepsilon$ and $N\ge N_\varepsilon$, statements (i) and (ii) of
Proposition~1 hold with $r_1=r$ and $C_1=C$ with probability at least
$1-\varepsilon$. In particular
$\hat\theta-\theta^*_T=O_p(T^4+N^{-1/2})$.}

\noindent\emph{Proof.} Let $\tilde x=(\tilde x_1,\dots,\tilde x_L)$ be the
coordinates in the fixed frame; they range over a compact product
$\mathcal S$ of the three sets. Let $\tilde J(\tilde x)$ be the Jacobian of
the model with respect to all $4L$ coordinates and $(c,d)$, at $c=d=0$. Its
entries are polynomials in $\tilde x$. The model is invariant under a
common spatial rotation of the points and of the tensors, so the three
infinitesimal rotations of the points lie in the kernel of $\tilde J$, and
$\tilde J$ has the same rank at $\tilde x$ and at the rotated configuration
$x$ that satisfies the gauge of (A3). When $x_1^1\neq0$ and $x_2^2\neq0$,
the three gauge conditions define a complement of the infinitesimal
rotations, and therefore $J(\xi,0,0)$ has rank $D$ if and only if
$\tilde J(\tilde x)$ has rank $4L+31$. The configurations of smaller rank
are the common zeros of the minors of order $4L+31$. These polynomials do
not all vanish identically: for five landmarks in each set this is the
exact computation quoted in Sect.~\ref{app:open}, and a further landmark
adds four columns, which are nonzero only in the rows of its own clocks,
where they equal $\pm2\eta(x_{\rm new}-x_j)$ and have rank four for a
suitable position, so that the rank increases by four. Hence the
configurations of smaller rank, together with those in which two points
coincide or the spatial parts of $\tilde x_1$ and $\tilde x_2$ are linearly
dependent, form a set $\mathcal Z$ of Lebesgue measure zero.

For $\mu>0$ let $\mathcal G_\mu\subset\mathcal S$ be the set on which the
singular value of order $4L+31$ of $\tilde J$, the distances between the
points, $|\vec{\tilde x}_1|$ and the area spanned by $\vec{\tilde x}_1$
and $\vec{\tilde x}_2$ are all at least $\mu$. These sets are compact,
they increase to $\mathcal S\setminus\mathcal Z$ as $\mu\to0$, and the
gauge rotation is continuous on them. Their images $\mathcal K_\mu$ in the
gauge of (A3) are compact and satisfy (A3) and (A4), and for small $T$ the
landmarks lie in $I(p,q)$ by Lemma~1. Since the density is bounded, the
probability that $\tilde x\notin\mathcal G_\mu$ is at most a constant times
the Lebesgue measure of $\mathcal S\setminus\mathcal G_\mu$, which tends to
zero as $\mu\to0$, uniformly in $T$. Given $\varepsilon$, choose $\mu$ such
that this probability is at most $\varepsilon/2$, and apply Proposition~1
and Lemma~4 with $\mathcal K=\mathcal K_\mu$, conditionally on the
landmarks. \hfill$\square$

\medskip
The constants of Corollary~2 depend on $\varepsilon$ through $\mu$ and are
not controlled; the remarks of Sect.~\ref{app:open} on the optimizer, the
coverage and the constants apply unchanged.

\subsection{Relation to the theorem of Madsen}
Theorem~4.9 and Corollary~4.10 of Ref.~\cite{madsen} hold for a smooth
globally hyperbolic spacetime of any dimension $d\ge2$, without field
equations. The input is a Poisson sprinkling of known density in a fixed
compact causal diamond. Anchors are selected from the sample by tests on
counts, and the existence of accepted anchors is proved. The output is the
full Riemann tensor at every sample point that passes the tests, in a
reconstructed frame, and the bound
$C[\log(e+\rho)/\rho]^{1/[2(3d+4)]}$ holds simultaneously at all these
points with probability at least $1-\rho^{-n}$ for every $n$, with a
constant that depends on $n$, and almost surely for large $\rho$.

Proposition~1 is weaker in each of these respects. It is restricted to
$d=4$ and to vacuum, the tips are given, the conditioning of the landmarks
and the choice of the local solution are assumed, the tensor is estimated
at one point, and the convergence is in probability. Under these
assumptions the exponent is $1/6$ instead of $1/32$. The two exponents do
not refer to the same problem, and the comparison does not show that one
construction is more efficient than the other. In particular, the rate
$\rho^{-1/6}$ relies on vacuum, which makes the volume clocks accurate to
relative order $T^4$ (Lemma~2); without it the clocks are accurate only to
relative order $T^2$ and the present estimator is not consistent.

\subsection{Numerical check of the orders}
The script \texttt{check\_convergence\_orders.py} in the repository checks
the orders used above on exact rotating plane waves
(Sect.~\ref{tests}), for a fixed configuration of 48 landmarks in windows
scaled by $1$, $\tfrac12$, $\tfrac14$ and $\tfrac18$ and with exact
interval volumes in place of counts. Over the three successive halvings of
the window, the residual of the fifth-degree expansion of Lemma~1(a)
decreases by factors of $15.1$, $15.8$ and $15.9$, and the relative volume
correction of Lemma~2 by $15.8$, $15.9$ and $16.0$, to be compared with
$16$. The relative error of the fitted tensor decreases by $4.3$, $4.1$
and $4.0$ without $\sigma_2$ and by $3.8$, $4.0$ and $4.0$ with it, to be
compared with $4$, and the error $|\hat\theta-\theta^*_T|$ of all
parameters by $22.8$, $18.0$ and $16.6$ without $\sigma_2$ and by $26.9$,
$17.4$ and $16.6$ with it, to be compared with $16$. The same script
verifies the statement on Myrheim's formula made after Lemma~2 and
performs the exact rank computation
quoted in Sect.~\ref{app:open}. This checks the orders in one family of
spacetimes and one configuration; it does not test the assumptions.
 in note.tex.
% Packages: amsmath and amssymb only (already loaded by note.tex); no amsthm.
% Labels of note.tex used: setting, estimator, tests, limits, clock, cov, model, square, ball.
% Citation keys used: brewin, gibbons, wang, madsen, and ONE new key, oneill:
%   \bibitem{oneill} O'Neill, B.: Semi-Riemannian Geometry with Applications to
%   Relativity. Academic Press, New York (1983)
% Check of the orders and of the rank certificate: python check_convergence_orders.py
\section{A conditional convergence statement}\label{app:conv}
This appendix states what we can prove about the fit of
Sect.~\ref{estimator} in shrinking windows, and what we cannot. The
statement is conditional: it assumes that the landmarks form a
well-conditioned configuration, and it concerns the solution of the fit that
lies near the true geometry. It does not show that the selection rule and
the optimizer of Sect.~\ref{estimator} meet these two conditions, and it
says nothing about the coverage of the ball \eqref{ball}. Only orders are
established; no constant is evaluated.

\subsection{Setting and assumptions}\label{app:setting}
We write $|\cdot|$ for the Euclidean norm of a vector of components,
$\|\cdot\|$ for the corresponding operator norm, and $A\succeq B$ if $A-B$
is positive semidefinite. For two events $x,x'$ of a convex normal
neighbourhood, $s(x,x')=g(v,v)$ with $v=\exp_x^{-1}x'$ is the signed squared
geodesic interval, so that $s=-\tau^2$ for timelike separation, and the
quantity $\sigma(x,y)$ of Sect.~\ref{estimator} is $s/T^2$.

\begin{description}
\item[(A1) Geometry.] $(M,g)$ is a time-oriented four-dimensional Lorentzian
manifold with smooth metric, $o\in M$, and $(e_a)$ is an orthonormal frame
at $o$ with $e_0$ future-directed. $U$ is a convex normal neighbourhood of
$o$, that is, an open set that is a normal neighbourhood of each of its
points; then any two points of $U$ are joined by a unique geodesic segment
in $U$, and $(x,x')\mapsto\exp_x^{-1}x'$ is smooth on $U\times U$
\cite[Chap.~5]{oneill}. The metric is Ricci-flat on $U$, and in the normal
coordinates based at $(o,e_a)$ its matrix of components satisfies
$\|g_{ab}(x)-\eta_{ab}\|\le1/10$ on $U$. Every point of a Ricci-flat region
has arbitrarily small neighbourhoods of this kind. For $0<T\le T_0$ the tips
are $p=\exp_o(-\tfrac T2e_0)$ and $q=\exp_o(\tfrac T2e_0)$, and, as in
Sect.~\ref{setting}, $I(p,q)=J^+(p)\cap J^-(q)\subset U$; this holds for
small $T$ if $M$ is strongly causal at $o$. The proofs use finitely many
derivatives, six of $s$ and four of the metric components in normal
coordinates; we have not determined the least differentiability of $g$ that
provides them.
\item[(A2) Sampling.] Conditionally on the landmarks, the $N$ measurement
events are independent and uniformly distributed on $I(p,q)$ with respect to
the volume of $g$. This is the case when the landmarks are taken from an
independent selection sample.
\item[(A3) Landmark configuration.] There are $L$ landmarks
$\ell_1,\dots,\ell_L\in I(p,q)$; put $\ell_0=p$ and $\ell_{L+1}=q$. Their
normal coordinates in units of $T$ are $x_i\in\mathbb R^4$, that is,
$\ell_i=\exp_o(Tx_i^ae_a)$, so that $x_0=(-\tfrac12,0,0,0)$ and
$x_{L+1}=(\tfrac12,0,0,0)$. The spatial axes are those of the gauge of
Sect.~\ref{estimator}: $x_1=(x_1^0,x_1^1,0,0)$ with $x_1^1>0$ and
$x_2=(x_2^0,x_2^1,x_2^2,0)$ with $x_2^2>0$. This fixes the frame up to the
reflection of $e_3$, and we fix one of the two frames. The remaining $4L-3$
coordinates form $\xi\in\mathbb R^{4L-3}$. The clocks are the pairs
$(\ell_i,\ell_j)$ with $(i,j)$ in a fixed set $\mathcal P$ of $m$ index
pairs, $(i,j)\neq(0,L+1)$. We assume that $\xi$ lies in a fixed compact set
$\mathcal K\subset\mathbb R^{4L-3}$ such that, for every configuration in
$\mathcal K$, the $L+2$ points $x_i$ are distinct, $x_1,\dots,x_L$ lie in
the open flat interval $\{|x^0|+|\vec x|<\tfrac12\}$, and $x_j-x_i$ is
future-directed timelike with respect to $\eta$ for every
$(i,j)\in\mathcal P$.
\item[(A4) Conditioning.] Let $\theta=(\xi,c,d)\in\mathbb R^D$, $D=4L+31$,
where $c\in\mathbb R^{10}$ and $d\in\mathbb R^{24}$ are the components of the
curvature and of its gradient in the bases of Sect.~\ref{estimator} (there
$d$ is called $g$, a letter needed here for the metric). Let
$\sigma(\theta)\in\mathbb R^m$ be the model \eqref{model} for the pairs of
$\mathcal P$, with $\sigma_2$ as in \eqref{square} or with $\sigma_2=0$, and
$J(\theta)=\partial\sigma/\partial\theta$. We assume that the $m\times D$
matrix $J(\xi,0,0)$ has rank $D$ for every $\xi\in\mathcal K$.
\item[(A5) Weights.] The fit minimizes
$\chi^2(\theta)=N[y-\sigma(\theta)]^{\!\top}\widehat W[y-\sigma(\theta)]$
with a symmetric matrix $\widehat W$ that may depend on the data, and there
are constants $0<w_-\le w_+$ such that $w_-I\preceq\widehat W\preceq w_+I$
with probability at least $1-\varpi_N$, where $\varpi_N\to0$ as
$N\to\infty$. For the weight of Sect.~\ref{estimator},
$\widehat W=(N\widehat\Sigma)^{-1}$, this is proved in
Lemma~4 below.
\end{description}

Assumptions (A3) and (A4) are substantive, and (A4) is a condition on the
shape of the configuration only. By compactness there are $\tau_0>0$ and
$\mu>0$ such that $\eta(x_j-x_i,x_j-x_i)\le-\tau_0^2$ for
$(i,j)\in\mathcal P$ and the smallest singular value of $J(\xi,0,0)$ is at
least $\mu$, for all $\xi\in\mathcal K$; all constants below depend on
these numbers. The spatial axes of the frame $(e_a)$ depend on the
landmarks; since the orthonormal frames at $o$ with the given $e_0$ form a
compact set, all estimates below are uniform over them. The true parameter
is $\theta^*_T=(\xi,c^*,d^*)$, where $\xi$
is the true configuration and $c^*$, $d^*$ are the components of
$T^2C_{abcd}(o)$ and of $T^3\nabla_eC_{abcd}(o)$ in the frame $(e_a)$. The
measured clocks are $y_{ij}=-(n_{ij}/N)^{1/2}$, with $n_{ij}$ the number of
measurement events in $I(\ell_i,\ell_j)$, and the continuum clocks are
$\bar y_{ij}=-[V(\ell_i,\ell_j)/V(p,q)]^{1/2}$.

\subsection{Statement}
\noindent\textbf{Proposition 1.} \emph{Assume (A1)--(A5). There are
constants $T_1,r_1>0$ and $C_1,C_2<\infty$, which depend only on the
geometry of $(U,g)$, on $\mathcal K$, $\mathcal P$ and on $w_\pm$, such that
for all $0<T\le T_1$, all $N\ge2$ and every landmark configuration allowed
by (A3) the following holds with probability at least $1-\varpi_N-C_2/N$.}
\begin{enumerate}
\item[(i)] \emph{$\chi^2$ has exactly one critical point
$\hat\theta=(\hat\xi,\hat c,\hat d)$ in the closed ball of radius $r_1$
around $\theta^*_T$, and $\hat\theta$ is the unique minimizer of $\chi^2$
on this ball.}
\item[(ii)] \emph{$|\hat\theta-\theta^*_T|\le C_1\left(T^4+|y-\bar y|\right)$.}
\end{enumerate}
\emph{Moreover $\operatorname{E}|y-\bar y|^2\le m/N$. Consequently, for
every $t>0$, with probability at least $1-\varpi_N-C_2/N-m/t^2$,}
\begin{equation}
 |\hat\theta-\theta^*_T|\le C_1\left(T^4+tN^{-1/2}\right),\qquad
 \left|\hat c/T^2-c^*/T^2\right|\le C_1\left(T^2+tN^{-1/2}T^{-2}\right),
 \label{conv:rate}
\end{equation}
\emph{where $c^*/T^2$ are the components of the Weyl tensor $C_{abcd}(o)$.
The same holds, with other constants, in the tidal norm of
Sect.~\ref{setting}.}

\medskip
In short, $\hat\theta-\theta^*_T=O(T^4)+O_p(N^{-1/2})$ in the normalized
variables of the fit. The statement and its proof are the same with and
without $\sigma_2$, because $\sigma_2=O(|c|^2)$; in particular the
coefficient of the volume term in \eqref{square} is not used. The estimator
returns $\hat c$, an estimate of $T^2C(o)$. The second inequality of
\eqref{conv:rate} refers to the tensor in physical units, which requires
$T$, or equivalently bounds the relative error when $C(o)\neq0$.

\medskip
\noindent\textbf{Corollary 1.} \emph{Let the measurement events be a
Poisson sprinkling of density $\rho$ in $I(p,q)$, with lengths measured in
a fixed unit, and let (A1)--(A5) hold. If $T\to0$ and $\rho T^8\to\infty$,
then the solution $\hat\theta$ of Proposition~1 exists with probability
tending to one, and}
\begin{equation}
 \hat c/T^2-C(o)=O(T^2)+O_p\!\left(\rho^{-1/2}T^{-4}\right)\to0
 \quad\text{in probability}.
\end{equation}
\emph{For $T$ proportional to $\rho^{-1/12}$ the error is
$O_p(\rho^{-1/6})$.}

\subsection{Proofs}
\noindent\textbf{Lemma 1} (normal coordinates). \emph{Let $\mathcal F$ be a
compact set of orthonormal frames $(x,f)$ with $x\in U$. Let
$w\in\mathbb R^4$ be the normal coordinates based at $(x,f)$, and let
$s_{x,f}(w,w')$ and $\gamma_{x,f}(w)$ be the function $s$ and the matrix of
metric components in these coordinates. There are $r_0>0$ and $C_0$ such
that for all $(x,f)\in\mathcal F$ and $|w|,|w'|\le r_0$, with
$\Delta=w-w'$:}
\begin{enumerate}
\item[(a)] \emph{$s_{x,f}(w,w')=\eta(\Delta,\Delta)-\tfrac13R(w,w',w,w')
-\tfrac1{12}(w+w')^e\nabla_eR(w,w',w,w')+\varrho$ with
$|\varrho|\le C_0(|w|+|w'|)^6$, where $R$ and $\nabla R$ are the components
of the curvature and its gradient at $x$ in the frame $f$;}
\item[(b)] \emph{$s_{x,f}(w,w')-\eta(\Delta,\Delta)$ is bounded by
$C_0(|w|+|w'|)^4$ and its first derivatives by $C_0(|w|+|w'|)^3$;}
\item[(c)] $\bigl|\,|\det\gamma_{x,f}(w)|^{1/2}-1\,\bigr|\le C_0|w|^4$;
\item[(d)] $s_{x,f}(0,w)=\eta(w,w)$.
\end{enumerate}

\noindent\emph{Proof.} The maps $(x,f,w,w')\mapsto s_{x,f}(w,w')$ and
$(x,f,w)\mapsto\gamma_{x,f}(w)$ are smooth on an open set that contains
$\mathcal F\times\{0\}\times\{0\}$, hence on
$\mathcal F\times\{|w|,|w'|\le r_0\}$ for some $r_0>0$. Statement (d) is the
definition of $s$. For (a) apply Taylor's theorem in $(w,w')$ at the
origin: the derivatives of order six are bounded on this compact set, and
the Taylor polynomial of degree five is the expansion of the squared
geodesic interval in normal coordinates, Sect.~9 of Ref.~\cite{brewin}
(Eqs.~(11.21a)--(11.21c) in arXiv:0903.2087v3). That polynomial has no terms
of degree three, which gives (b) in the same way. For (c),
$\gamma_{ab}(w)=\eta_{ab}-\tfrac13R_{acbd}w^cw^d
-\tfrac16\nabla_eR_{acbd}w^cw^dw^e+O(|w|^4)$ \cite{brewin}, and the traces
of the second and third terms with $\eta^{ab}$ contain only the Ricci tensor
and its gradient, which vanish; hence $\det\gamma=-1+O(|w|^4)$, with a
remainder bounded as before. \hfill$\square$

\medskip
\noindent\textbf{Lemma 2} (volume of a small vacuum interval). \emph{Let
$\mathcal T$ be a compact set of future-directed unit timelike vectors
$(x,u)$ with $x\in U$. There are $\tau_1>0$ and $C_V$ such that, whenever
$(x,u)\in\mathcal T$, $0<\tau\le\tau_1$, $x'=\exp_x(\tau u)$ and
$J^+(x)\cap J^-(x')\subset U$,}
\begin{equation}
 \left|V(x,x')-\tfrac{\pi}{24}\tau^4\right|\le C_V\tau^8 .
 \label{conv:volume}
\end{equation}

\noindent\emph{Proof.} Complete $u$ to an orthonormal frame $f$ with
$f_0=u$. The frames so obtained form a compact set, to which Lemma~1
applies. We use the normal coordinates $w$ based at $(x,f)$, in which $x'$
is $\tau\hat e$ with $\hat e=(1,0,0,0)$, write $s=s_{x,f}$, and put
$\Xi(w)=s(0,w)-s(w,\tau\hat e)$. We show that the half
$H_-=\{z\in J^+(x)\cap J^-(x'):\Xi\ge0\}$ of the interval has volume
$\tfrac{\pi}{48}\tau^4+O(\tau^8)$. The other half, where $\Xi\le0$, is
treated in the same way from $x'$ with the time orientation reversed, and
the common boundary $\Xi=0$ has measure zero. All $O(\cdot)$ below are
uniform over $\mathcal T$ and over the directions $\nu$.

(1) \emph{Localization.} By (A1), every causal vector $v$ at a point of $U$
has $|\vec v|\le\tfrac54|v^0|$ in the coordinates of (A1), with $v^0>0$ if
it is future-directed. Hence $x^0$ increases along future-directed causal
curves in $U$, and the coordinate differences satisfy
$|z-x|\le\tfrac94|x'-x|$ for every $z\in J^+(x)\cap J^-(x')$. Since the
change to the coordinates $w$ is smooth, uniformly over the frames, there
is $C\ge1$ such that $|w|\le C\tau$ for all these $z$.

(2) \emph{Rays.} In a convex normal neighbourhood, $z\in J^+(x)$ if and
only if $\exp_x^{-1}z$ is zero or future-directed causal
\cite[Lemma~14.2]{oneill}, that is, if and only if $w=t\nu$ with $t\ge0$
and $\nu\in\mathcal N=\{(1,\vec\nu):|\vec\nu|\le1\}$. For $\nu\in\mathcal N$
let $\phi_\nu(t)=\Xi(t\nu)$. By Lemma~1(d),
$\phi_\nu(t)=t^2\eta(\nu,\nu)-s(t\nu,\tau\hat e)$, and
$\eta(t\nu-\tau\hat e,t\nu-\tau\hat e)=t^2\eta(\nu,\nu)+2\tau t-\tau^2$.
By Lemma~1(b), for $0\le t\le C\tau$,
\begin{equation}
 \phi_\nu(t)=\tau^2-2\tau t+O(\tau^4),\qquad
 \phi_\nu'(t)=-2\tau+O(\tau^3).
\end{equation}
For $\tau\le\tau_1$ small, $\phi_\nu$ is therefore strictly decreasing on
$[0,C\tau]$, with $\phi_\nu(0)=\tau^2$ and $\phi_\nu(\tau)<0$. It has a
single zero $t_*(\nu)\in(0,\tau)$, which depends continuously on $\nu$.

(3) \emph{The half-interval is
$\{t\nu:\nu\in\mathcal N,\ 0\le t\le t_*(\nu)\}$.} If $z\in H_-$, then
$w=t\nu$ with $t\le C\tau$ by (1) and (2), and $\phi_\nu(t)\ge0$ gives
$t\le t_*(\nu)$. Conversely, let $0\le t\le t_*(\nu)$. On the segment from
$0$ to $t\nu$ we have $\Xi\ge0$, that is, $s(\cdot,x')\le s(x,\cdot)\le0$,
so the segment lies in the set of events causally related to $x'$, and it
does not contain $x'$, where $\Xi=-\tau^2$. The sets
$J^-(x')\setminus\{x'\}$ and $J^+(x')\setminus\{x'\}$ are disjoint and
closed in $U\setminus\{x'\}$, and the segment starts at $x\in J^-(x')$;
being connected, it lies in $J^-(x')$. Hence $t\nu\in H_-$.

(4) \emph{Volume.} With $w=(t,t\vec\nu)$ we have
$\mathrm d^4w=t^3\,\mathrm dt\,\mathrm d^3\vec\nu$, and Lemma~1(c) gives
\begin{equation}
 \operatorname{Vol}(H_-)=\int_{|\vec\nu|\le1}\int_0^{t_*(\nu)}
 \left[1+O(\tau^4)\right]t^3\,\mathrm dt\,\mathrm d^3\vec\nu
 =\frac14\int_{|\vec\nu|\le1}t_*(\nu)^4\,\mathrm d^3\vec\nu+O(\tau^8).
 \label{conv:half}
\end{equation}

(5) \emph{Expansion of $t_*$.} Insert Lemma~1(a) with $w=t\nu$ and
$w'=\tau\hat e$. Since
$R(t\nu,\tau\hat e,t\nu,\tau\hat e)=t^2\tau^2E_{ij}\nu^i\nu^j$ with
$E_{ij}=R_{0i0j}$,
\begin{equation}
 \phi_\nu(t)=\tau^2-2\tau t+\tfrac13t^2\tau^2E_{ij}\nu^i\nu^j
 +\tfrac1{12}t^2\tau^2\left[(t+\tau)\nabla_0E_{ij}
 +t\,\nu^k\nabla_kE_{ij}\right]\nu^i\nu^j+O(\tau^6).
\end{equation}
The zero satisfies $t_*=\tfrac\tau2+O(\tau^3)$, and then
\begin{equation}
 t_*(\nu)=\frac\tau2\left[1+\frac{\tau^2}{12}E_{ij}\nu^i\nu^j
 +\frac{\tau^3}{96}\left(3\nabla_0E_{ij}+\nu^k\nabla_kE_{ij}\right)
 \nu^i\nu^j+O(\tau^4)\right].
\end{equation}
The integral of $\nu^i\nu^j$ over the unit ball is
$\tfrac{4\pi}{15}\delta^{ij}$, and that of $\nu^i\nu^j\nu^k$ vanishes. In
vacuum $E_{ii}=R_{00}=0$ and $\nabla_0E_{ii}=\nabla_0R_{00}=0$, so the terms
of relative order $\tau^2$ and $\tau^3$ integrate to zero, and
\eqref{conv:half} gives
$\operatorname{Vol}(H_-)=\tfrac14\cdot\tfrac{4\pi}3(\tau/2)^4+O(\tau^8)
=\tfrac{\pi}{48}\tau^4+O(\tau^8)$. \hfill$\square$

\medskip
Lemma~2 uses the field equations only through $R_{ab}=0$, and it gives only
the order of the correction. As a check of the method, without that
condition the same computation, with the volume element
$1-\tfrac16R_{ab}w^aw^b$, gives
$V=\tfrac{\pi}{24}\tau^4[1-\tfrac1{180}R^a{}_a\tau^2+\tfrac1{30}R_{00}\tau^2
+O(\tau^3)]$, which is the formula of Myrheim, Eqs.~(7), (14) and (23) of
Ref.~\cite{gibbons}. In vacuum the first correction is of relative order
$\tau^4$, in agreement with Eq.~(79) of Ref.~\cite{wang}; its coefficient, on
which \eqref{square} and Ref.~\cite{wang} differ, plays no role here.

\medskip
\noindent\textbf{Lemma 3} (clocks). \emph{Assume (A1)--(A3). (a) There are
$T_1>0$ and $C_b$ such that $|\bar y_{ij}-\sigma_{ij}(\theta^*_T)|\le
C_bT^4$ for all $T\le T_1$, $\xi\in\mathcal K$ and $(i,j)\in\mathcal P$.
(b) $\operatorname{E}|y-\bar y|^2\le m/N$.}

\noindent\emph{Proof.} (a) Lemma~1(a) at $(o,e)$ with $w=Tx_i$ and
$w'=Tx_j$, and $R_{abcd}=C_{abcd}$ in vacuum, give
\begin{equation}
 \frac{s(\ell_i,\ell_j)}{T^2}=\eta(\Delta,\Delta)
 -\frac{T^2}3C(x_i,x_j,x_i,x_j)
 -\frac{T^3}{12}(x_i+x_j)^e\nabla_eC(x_i,x_j,x_i,x_j)+O(T^4),
\end{equation}
with $\Delta=x_i-x_j$ and the tensors at $o$. The right-hand side is the
model \eqref{model} at $\theta^*_T$ without $\sigma_2$, up to $O(T^4)$, and
$\sigma_2(\theta^*_T)=O(|c^*|^2)=O(T^4)$. Since
$\eta(\Delta,\Delta)\le-\tau_0^2$, for small $T$ we have
$\tau_{ij}^2=-s(\ell_i,\ell_j)\ge\tau_0^2T^2/2$ and $\tau_{ij}=O(T)$. In
the coordinates of (A1) the vector $\exp_{\ell_i}^{-1}\ell_j$ has components
$T(x_j-x_i)+O(T^2)$, so for small $T$ it is future-directed and $\ell_j$
lies in the timelike future of $\ell_i$. The unit tangent at $\ell_i$ of the
geodesic to $\ell_j$ then has bounded components, and these vectors lie in a
compact set $\mathcal T$.
Moreover $J^+(\ell_i)\cap J^-(\ell_j)\subset I(p,q)\subset U$. Lemma~2
therefore gives $V(\ell_i,\ell_j)=\tfrac{\pi}{24}\tau_{ij}^4[1+O(T^4)]$, and
likewise $V(p,q)=\tfrac{\pi}{24}T^4[1+O(T^4)]$. Hence
$\bar y_{ij}=-(\tau_{ij}^2/T^2)[1+O(T^4)]=s(\ell_i,\ell_j)/T^2+O(T^4)$.

(b) Given the landmarks, $n_{ij}$ is binomial with parameters $N$ and
$\pi_{ij}=V(\ell_i,\ell_j)/V(p,q)=\bar y_{ij}^2$. Since
$|\sqrt a-\sqrt b|\le|a-b|/\sqrt b$,
\begin{equation}
 |y_{ij}-\bar y_{ij}|\le\frac{|n_{ij}/N-\pi_{ij}|}{\sqrt{\pi_{ij}}},
 \qquad\operatorname{E}|y_{ij}-\bar y_{ij}|^2\le\frac{1-\pi_{ij}}N .
\end{equation}
Summing over the $m$ pairs gives (b). The clocks are strongly correlated,
see \eqref{cov}, but no independence is used. \hfill$\square$

\medskip
\noindent\textbf{Lemma 4} (weights). \emph{Let $N\widehat\Sigma=\widehat
D\,\widehat\Gamma\,\widehat D$ be the covariance of Sect.~\ref{estimator},
where $\widehat\Gamma$ is the empirical covariance matrix of the membership
indicators of the $m$ intervals and
$\widehat D=\operatorname{diag}(\tfrac12\hat\pi_{ij}^{-1/2})$ with
$\hat\pi_{ij}=n_{ij}/N$. Under (A1)--(A3) there are $T_2>0$,
$0<w_-\le w_+$ and $C_W$ such that $\widehat W=(N\widehat\Sigma)^{-1}$
satisfies (A5) with $\varpi_N=C_W/N$ for all $T\le T_2$.}

\noindent\emph{Proof.} Let $\Gamma_T(\xi)$ be the covariance matrix of the
indicators under the sampling law, and $\Gamma_0(\xi)$ the same matrix for
the flat metric $\eta$ and the points $x_i$.

(1) \emph{$\Gamma_0(\xi)$ is positive definite.} Otherwise a combination
$\sum\lambda_{ij}1_{ij}$ of the indicators with $\lambda\neq0$ is constant
almost everywhere in the flat interval. Fix a pair $(i,j)$. The edge of its
interval, where the future cone of $x_i$ meets the past cone of $x_j$, is a
round two-sphere, and the only null cones that contain it are those of
$x_i$ and $x_j$; the null cone of any other point meets it in the zero set
of a quadratic polynomial that does not vanish identically on the sphere,
a closed set without interior points. Choose $z$ on the edge and on none of
the other cones, and a ball $B$ around $z$ that meets no null cone of a
point $x_k$ other than the future cone of $x_i$ and the past cone of $x_j$.
In $B$, every indicator other than $1_{ij}$ is a constant, or a constant
times the indicator of $J^+(x_i)$ or of $J^-(x_j)$, while $1_{ij}$ is the
product of these two. The two cones cross transversally at $z$. If
$1\le i,j\le L$, all four combinations of the two indicators occur on sets
of positive measure in $B$, and constancy of the sum forces
$\lambda_{ij}=0$. If $i=0$, only $B\cap J^+(p)$ carries probability, and no
interval $I(x_k,x_j)$ with $k\neq0$ meets $B$, because
$z\in\partial J^+(p)$ whereas $x_k$ lies in the open interval; the two
values of the indicator of $J^-(x_j)$ again force $\lambda_{0j}=0$. The
case $j=L+1$ is symmetric.

(2) In the coordinates $x/T$ the metric has components $g_{ab}(Tx)$. These,
the inverse exponential map and the volume density converge to their flat
values as $T\to0$, uniformly on compact sets. The indicator of each
interval is therefore continuous in $(T,\xi)$ and in the spatial frame at
$o$ at every point off the boundary of the limiting interval, and by
dominated convergence $\Gamma_T(\xi)$ is continuous in these variables on
$[0,T_0]\times\mathcal K$ times the compact set of frames, with the value
$\Gamma_0(\xi)$ at $T=0$. With (1) and compactness,
$\Gamma_T(\xi)\succeq4\gamma I$ for some $\gamma>0$, all $\xi\in\mathcal K$,
all frames and all $T\le T_2$. In the same way $\pi_{ij}\ge\pi_->0$.

(3) $\widehat\Gamma$ and $\hat\pi$ are built from averages of bounded
variables, so $\operatorname{E}\|\widehat\Gamma-\Gamma_T\|^2\le Cm^2/N$ and
$\operatorname{E}|\hat\pi-\pi|^2\le m/N$. By Chebyshev's inequality the
event that $\|\widehat\Gamma-\Gamma_T\|\le2\gamma$ and
$\hat\pi_{ij}\ge\pi_-/2$ for all pairs has probability at least $1-C_W/N$.
On this event the diagonal entries of $\widehat D$ lie between $\tfrac12$
and $(2\pi_-)^{-1/2}$, $\widehat\Gamma\succeq2\gamma I$ and
$\operatorname{tr}\widehat\Gamma\le m/2$, hence
$\tfrac\gamma2I\preceq N\widehat\Sigma\preceq\tfrac{m}{4\pi_-}I$.
\hfill$\square$

\medskip
\noindent\textbf{Lemma 5} (local minimizer). \emph{Let $F$ be a twice
continuously differentiable map from the closed ball of radius $r_0$ around
$\theta^*\in\mathbb R^D$ to $\mathbb R^m$, with $\|DF\|\le M_1$ and
$\|D^2F\|\le M_2$. Let $W$ be symmetric with $0\preceq W\preceq w_+I$ and
$DF(\theta^*)^{\!\top}W\,DF(\theta^*)\succeq\lambda I$, $\lambda>0$. For
$y\in\mathbb R^m$ let $Q(\theta)=\tfrac12[y-F(\theta)]^{\!\top}
W[y-F(\theta)]$ and $\epsilon=|y-F(\theta^*)|$, and put}
\begin{equation}
 r_1=\min\left\{r_0,\frac{\lambda}{12M_1M_2w_+}\right\},\qquad
 \epsilon_1=\min\left\{\frac{\lambda}{4M_2w_+},
 \frac{\lambda r_1}{8M_1w_+}\right\}.
\end{equation}
\emph{If $\epsilon\le\epsilon_1$, then $Q$ has exactly one critical point
$\hat\theta$ in the closed ball of radius $r_1$ around $\theta^*$, it is the
unique minimizer of $Q$ on this ball, and
$|\hat\theta-\theta^*|\le4M_1w_+\epsilon/\lambda$.}

\noindent\emph{Proof.} Let $|\theta-\theta^*|\le r_1$, let $v$ be a unit
vector and $\rho=y-F(\theta)$. Then
\begin{equation}
 v^{\!\top}D^2Q(\theta)\,v=[DF(\theta)v]^{\!\top}W[DF(\theta)v]
 -\rho^{\!\top}W\,D^2F(\theta)[v,v].
\end{equation}
Here $|\rho|\le\epsilon+M_1r_1$ and $\|DF(\theta)-DF(\theta^*)\|\le M_2r_1$,
so that
$v^{\!\top}D^2Q(\theta)v\ge\lambda-2M_1M_2w_+r_1-M_2w_+(\epsilon+M_1r_1)
\ge\lambda/2$. Thus $Q$ is strongly convex on the ball. It has a unique
minimizer $\hat\theta$ there and at most one critical point, and
$Q(\theta)\ge Q(\theta^*)+DQ(\theta^*)(\theta-\theta^*)
+\tfrac\lambda4|\theta-\theta^*|^2$. Since
$|DQ(\theta^*)|=|DF(\theta^*)^{\!\top}W[y-F(\theta^*)]|\le M_1w_+\epsilon$,
we get $Q(\theta)>Q(\theta^*)$ whenever
$|\theta-\theta^*|>4M_1w_+\epsilon/\lambda$. Hence
$|\hat\theta-\theta^*|\le4M_1w_+\epsilon/\lambda\le r_1/2$, so $\hat\theta$
is an interior point and therefore a critical point. \hfill$\square$

\medskip
\noindent\emph{Proof of Proposition 1.} The map $\sigma$ is a polynomial in
$\theta$. Let $M_1$ and $M_2$ bound its first and second derivatives on the
set of all $\theta$ at distance at most one from
$\{(\xi,c,d):\xi\in\mathcal K,\ |c|,|d|\le1\}$; for small $T$ this set
contains the ball of radius $r_0=1$ around $\theta^*_T$. Because $\sigma_2$
is quadratic in $c$, the matrix $J(\xi,0,0)$ does not depend on it, and
$\|J(\theta^*_T)-J(\xi,0,0)\|\le M_2(|c^*|+|d^*|)=O(T^2)$. By (A4) the
smallest singular value of $J(\theta^*_T)$ is at least $\mu/2$ for
$T\le T_1$, so that on the event of (A5)
$J(\theta^*_T)^{\!\top}\widehat WJ(\theta^*_T)\succeq\lambda I$ with
$\lambda=w_-\mu^2/4$. Apply Lemma~5 on this event with $F=\sigma$,
$W=\widehat W$ and $\theta^*=\theta^*_T$; the functions $\chi^2$ and $Q$
differ by the factor $2N$. By Lemma~3(a),
$\epsilon=|y-\sigma(\theta^*_T)|\le|y-\bar y|+\sqrt m\,C_bT^4$. Reduce $T_1$
so that $\sqrt m\,C_bT_1^4\le\epsilon_1/2$. By Lemma~3(b) and Chebyshev's
inequality, $|y-\bar y|\le\epsilon_1/2$ except on an event of probability at
most $C_2/N$ with $C_2=4m/\epsilon_1^2$. Lemma~5 then gives (i), and (ii)
with $C_1=4M_1w_+\max(1,\sqrt m\,C_b)/\lambda$. Finally
$|y-\bar y|\le tN^{-1/2}$ except with probability $m/t^2$, which gives
\eqref{conv:rate}, the second inequality because
$|\hat c-c^*|\le|\hat\theta-\theta^*_T|$. \hfill$\square$

\medskip
\noindent\emph{Proof of Corollary 1.} Given their number $N$, the events of
a Poisson sprinkling in $I(p,q)$ satisfy (A2), and $N$ is Poisson with mean
$\rho V(p,q)=\tfrac{\pi}{24}\rho T^4[1+O(T^4)]$ by Lemma~2. If
$\rho T^4\to\infty$, then $N\ge\rho T^4/20$ with probability tending to
one, and on this event $N^{-1/2}T^{-2}\le\sqrt{20}\,\rho^{-1/2}T^{-4}$ and
$\varpi_N+C_2/N\to0$. Insert this into \eqref{conv:rate}. The two terms
balance for $T^6$ proportional to $\rho^{-1/2}$, and then
$\rho T^8\to\infty$. \hfill$\square$

\subsection{What is not proved}\label{app:open}
\emph{The selection rule.} Assumptions (A3) and (A4) are not shown to hold
for the landmarks produced by the rule of Sect.~\ref{estimator}, which are
events of a selection sample chosen by counts, with a clock set chosen by
counts. Assumption (A4) is not vacuous: for an explicit rational
configuration with five landmarks in each of three time bands, and with
the clocks between different bands and to the tips, exact modular
arithmetic gives $J(\xi,0,0)$ the full rank $D=91$. Since full rank is the
non-vanishing of a polynomial, (A4) then holds for almost every
configuration with these clocks (see Corollary~2). This identifies the
parameters; it does not bound $\mu$, on which the constants depend. In
that configuration the singular values of $J(\xi,0,0)$ range over eight
orders of magnitude.

\emph{The optimizer.} Proposition~1 concerns the critical point near
$\theta^*_T$. It does not show that the Levenberg--Marquardt iteration from
the four starts of Sect.~\ref{estimator} converges to this point, and it
does not exclude other local minima. The reflections of
Sect.~\ref{setting} map $\hat\theta$ to further minimizers with the same
$\chi^2$; these are the mirror images accounted for there. No global
identifiability is proved.

\emph{The coverage.} Nothing is proved about the ball \eqref{ball}. For
$T$ proportional to $\rho^{-1/12}$ the bias and the noise in
\eqref{conv:rate} are of the same order, so a ball built from the noise
covariance alone cannot be expected to cover, and the allowance $\delta$
remains an assumption. In the regime $N^{1/2}T^4\to0$, where the noise
dominates, one may expect $\hat\theta$ to be asymptotically normal with the
covariance used in Sect.~\ref{estimator}, but we have not proved this, and
no asymptotic statement bears on the coverage at the values of $N$ and $T$
of the tests.

\emph{The constants.} $C_1$ and $T_1$ depend on $\mu$ and on bounds for the
derivatives of the metric on $U$ up to the order that controls the sixth
derivatives of $s$. They are not computed, so Proposition~1 gives no bias
bound at a given window, such as the allowance of 2\% in \eqref{ball}. It
is not uniform over spacetimes without a bound on these derivatives, in
agreement with the obstruction of Sect.~\ref{limits}.

\emph{Further restrictions.} The tips are given events and are not chosen
from the causal set. Vacuum is essential: Ricci terms change the clocks
\eqref{clock} at relative order $T^2$, the order of the signal. The
implementation discards eigenvalues of the clock correlation matrix below
$10^{-8}$ times the largest. Lemma~4 covers this only when no eigenvalue is
discarded; otherwise the lower bound in (A5) has to be replaced by a lower
bound on $J^{\!\top}\widehat WJ$, which is all that the proof uses. The rate
is an upper bound for this estimator under these assumptions, and we claim
no optimality.

\subsection{An idealized random design}\label{app:random}
The conditioning assumption can be removed, in probability, for a
simplified design in which the landmarks are placed at random with a
bounded density. This is not the order-based rule of
Sect.~\ref{estimator}.

\medskip
\noindent\textbf{Corollary 2.} \emph{Let $B_-,B_0,B_+$ be compact subsets
of the open flat interval $\{|x^0|+|\vec x|<\tfrac12\}$ such that $x'-x$ is
future-directed timelike with respect to $\eta$ whenever
$x\in B_-,x'\in B_0$, or $x\in B_0,x'\in B_+$, or $x\in B_-,x'\in B_+$.
Assume (A1) and (A2), and let the landmarks be random, with at least five
in each of the three sets $\exp_o(TB_-)$, $\exp_o(TB_0)$, $\exp_o(TB_+)$,
their numbers being fixed. Assume that their normal coordinates in units
of $T$, with respect to a fixed frame at $o$, have a joint density bounded
by a constant independent of $T$. Let the clocks be all pairs of landmarks
from different sets and all pairs of a landmark and a tip, and let
$\widehat W$ be the weight of Lemma~4 or the identity. Then for every
$\varepsilon>0$ there are $r,C,T_\varepsilon,N_\varepsilon$ such that, for
$T\le T_\varepsilon$ and $N\ge N_\varepsilon$, statements (i) and (ii) of
Proposition~1 hold with $r_1=r$ and $C_1=C$ with probability at least
$1-\varepsilon$. In particular
$\hat\theta-\theta^*_T=O_p(T^4+N^{-1/2})$.}

\noindent\emph{Proof.} Let $\tilde x=(\tilde x_1,\dots,\tilde x_L)$ be the
coordinates in the fixed frame; they range over a compact product
$\mathcal S$ of the three sets. Let $\tilde J(\tilde x)$ be the Jacobian of
the model with respect to all $4L$ coordinates and $(c,d)$, at $c=d=0$. Its
entries are polynomials in $\tilde x$. The model is invariant under a
common spatial rotation of the points and of the tensors, so the three
infinitesimal rotations of the points lie in the kernel of $\tilde J$, and
$\tilde J$ has the same rank at $\tilde x$ and at the rotated configuration
$x$ that satisfies the gauge of (A3). When $x_1^1\neq0$ and $x_2^2\neq0$,
the three gauge conditions define a complement of the infinitesimal
rotations, and therefore $J(\xi,0,0)$ has rank $D$ if and only if
$\tilde J(\tilde x)$ has rank $4L+31$. The configurations of smaller rank
are the common zeros of the minors of order $4L+31$. These polynomials do
not all vanish identically: for five landmarks in each set this is the
exact computation quoted in Sect.~\ref{app:open}, and a further landmark
adds four columns, which are nonzero only in the rows of its own clocks,
where they equal $\pm2\eta(x_{\rm new}-x_j)$ and have rank four for a
suitable position, so that the rank increases by four. Hence the
configurations of smaller rank, together with those in which two points
coincide or the spatial parts of $\tilde x_1$ and $\tilde x_2$ are linearly
dependent, form a set $\mathcal Z$ of Lebesgue measure zero.

For $\mu>0$ let $\mathcal G_\mu\subset\mathcal S$ be the set on which the
singular value of order $4L+31$ of $\tilde J$, the distances between the
points, $|\vec{\tilde x}_1|$ and the area spanned by $\vec{\tilde x}_1$
and $\vec{\tilde x}_2$ are all at least $\mu$. These sets are compact,
they increase to $\mathcal S\setminus\mathcal Z$ as $\mu\to0$, and the
gauge rotation is continuous on them. Their images $\mathcal K_\mu$ in the
gauge of (A3) are compact and satisfy (A3) and (A4), and for small $T$ the
landmarks lie in $I(p,q)$ by Lemma~1. Since the density is bounded, the
probability that $\tilde x\notin\mathcal G_\mu$ is at most a constant times
the Lebesgue measure of $\mathcal S\setminus\mathcal G_\mu$, which tends to
zero as $\mu\to0$, uniformly in $T$. Given $\varepsilon$, choose $\mu$ such
that this probability is at most $\varepsilon/2$, and apply Proposition~1
and Lemma~4 with $\mathcal K=\mathcal K_\mu$, conditionally on the
landmarks. \hfill$\square$

\medskip
The constants of Corollary~2 depend on $\varepsilon$ through $\mu$ and are
not controlled; the remarks of Sect.~\ref{app:open} on the optimizer, the
coverage and the constants apply unchanged.

\subsection{Relation to the theorem of Madsen}
Theorem~4.9 and Corollary~4.10 of Ref.~\cite{madsen} hold for a smooth
globally hyperbolic spacetime of any dimension $d\ge2$, without field
equations. The input is a Poisson sprinkling of known density in a fixed
compact causal diamond. Anchors are selected from the sample by tests on
counts, and the existence of accepted anchors is proved. The output is the
full Riemann tensor at every sample point that passes the tests, in a
reconstructed frame, and the bound
$C[\log(e+\rho)/\rho]^{1/[2(3d+4)]}$ holds simultaneously at all these
points with probability at least $1-\rho^{-n}$ for every $n$, with a
constant that depends on $n$, and almost surely for large $\rho$.

Proposition~1 is weaker in each of these respects. It is restricted to
$d=4$ and to vacuum, the tips are given, the conditioning of the landmarks
and the choice of the local solution are assumed, the tensor is estimated
at one point, and the convergence is in probability. Under these
assumptions the exponent is $1/6$ instead of $1/32$. The two exponents do
not refer to the same problem, and the comparison does not show that one
construction is more efficient than the other. In particular, the rate
$\rho^{-1/6}$ relies on vacuum, which makes the volume clocks accurate to
relative order $T^4$ (Lemma~2); without it the clocks are accurate only to
relative order $T^2$ and the present estimator is not consistent.

\subsection{Numerical check of the orders}
The script \texttt{check\_convergence\_orders.py} in the repository checks
the orders used above on exact rotating plane waves
(Sect.~\ref{tests}), for a fixed configuration of 48 landmarks in windows
scaled by $1$, $\tfrac12$, $\tfrac14$ and $\tfrac18$ and with exact
interval volumes in place of counts. Over the three successive halvings of
the window, the residual of the fifth-degree expansion of Lemma~1(a)
decreases by factors of $15.1$, $15.8$ and $15.9$, and the relative volume
correction of Lemma~2 by $15.8$, $15.9$ and $16.0$, to be compared with
$16$. The relative error of the fitted tensor decreases by $4.3$, $4.1$
and $4.0$ without $\sigma_2$ and by $3.8$, $4.0$ and $4.0$ with it, to be
compared with $4$, and the error $|\hat\theta-\theta^*_T|$ of all
parameters by $22.8$, $18.0$ and $16.6$ without $\sigma_2$ and by $26.9$,
$17.4$ and $16.6$ with it, to be compared with $16$. The same script
verifies the statement on Myrheim's formula made after Lemma~2 and
performs the exact rank computation
quoted in Sect.~\ref{app:open}. This checks the orders in one family of
spacetimes and one configuration; it does not test the assumptions.
 in note.tex.
% Packages: amsmath and amssymb only (already loaded by note.tex); no amsthm.
% Labels of note.tex used: setting, estimator, tests, limits, clock, cov, model, square, ball.
% Citation keys used: brewin, gibbons, wang, madsen, and ONE new key, oneill:
%   \bibitem{oneill} O'Neill, B.: Semi-Riemannian Geometry with Applications to
%   Relativity. Academic Press, New York (1983)
% Check of the orders and of the rank certificate: python check_convergence_orders.py
\section{A conditional convergence statement}\label{app:conv}
This appendix states what we can prove about the fit of
Sect.~\ref{estimator} in shrinking windows, and what we cannot. The
statement is conditional: it assumes that the landmarks form a
well-conditioned configuration, and it concerns the solution of the fit that
lies near the true geometry. It does not show that the selection rule and
the optimizer of Sect.~\ref{estimator} meet these two conditions, and it
says nothing about the coverage of the ball \eqref{ball}. Only orders are
established; no constant is evaluated.

\subsection{Setting and assumptions}\label{app:setting}
We write $|\cdot|$ for the Euclidean norm of a vector of components,
$\|\cdot\|$ for the corresponding operator norm, and $A\succeq B$ if $A-B$
is positive semidefinite. For two events $x,x'$ of a convex normal
neighbourhood, $s(x,x')=g(v,v)$ with $v=\exp_x^{-1}x'$ is the signed squared
geodesic interval, so that $s=-\tau^2$ for timelike separation, and the
quantity $\sigma(x,y)$ of Sect.~\ref{estimator} is $s/T^2$.

\begin{description}
\item[(A1) Geometry.] $(M,g)$ is a time-oriented four-dimensional Lorentzian
manifold with smooth metric, $o\in M$, and $(e_a)$ is an orthonormal frame
at $o$ with $e_0$ future-directed. $U$ is a convex normal neighbourhood of
$o$, that is, an open set that is a normal neighbourhood of each of its
points; then any two points of $U$ are joined by a unique geodesic segment
in $U$, and $(x,x')\mapsto\exp_x^{-1}x'$ is smooth on $U\times U$
\cite[Chap.~5]{oneill}. The metric is Ricci-flat on $U$, and in the normal
coordinates based at $(o,e_a)$ its matrix of components satisfies
$\|g_{ab}(x)-\eta_{ab}\|\le1/10$ on $U$. Every point of a Ricci-flat region
has arbitrarily small neighbourhoods of this kind. For $0<T\le T_0$ the tips
are $p=\exp_o(-\tfrac T2e_0)$ and $q=\exp_o(\tfrac T2e_0)$, and, as in
Sect.~\ref{setting}, $I(p,q)=J^+(p)\cap J^-(q)\subset U$; this holds for
small $T$ if $M$ is strongly causal at $o$. The proofs use finitely many
derivatives, six of $s$ and four of the metric components in normal
coordinates; we have not determined the least differentiability of $g$ that
provides them.
\item[(A2) Sampling.] Conditionally on the landmarks, the $N$ measurement
events are independent and uniformly distributed on $I(p,q)$ with respect to
the volume of $g$. This is the case when the landmarks are taken from an
independent selection sample.
\item[(A3) Landmark configuration.] There are $L$ landmarks
$\ell_1,\dots,\ell_L\in I(p,q)$; put $\ell_0=p$ and $\ell_{L+1}=q$. Their
normal coordinates in units of $T$ are $x_i\in\mathbb R^4$, that is,
$\ell_i=\exp_o(Tx_i^ae_a)$, so that $x_0=(-\tfrac12,0,0,0)$ and
$x_{L+1}=(\tfrac12,0,0,0)$. The spatial axes are those of the gauge of
Sect.~\ref{estimator}: $x_1=(x_1^0,x_1^1,0,0)$ with $x_1^1>0$ and
$x_2=(x_2^0,x_2^1,x_2^2,0)$ with $x_2^2>0$. This fixes the frame up to the
reflection of $e_3$, and we fix one of the two frames. The remaining $4L-3$
coordinates form $\xi\in\mathbb R^{4L-3}$. The clocks are the pairs
$(\ell_i,\ell_j)$ with $(i,j)$ in a fixed set $\mathcal P$ of $m$ index
pairs, $(i,j)\neq(0,L+1)$. We assume that $\xi$ lies in a fixed compact set
$\mathcal K\subset\mathbb R^{4L-3}$ such that, for every configuration in
$\mathcal K$, the $L+2$ points $x_i$ are distinct, $x_1,\dots,x_L$ lie in
the open flat interval $\{|x^0|+|\vec x|<\tfrac12\}$, and $x_j-x_i$ is
future-directed timelike with respect to $\eta$ for every
$(i,j)\in\mathcal P$.
\item[(A4) Conditioning.] Let $\theta=(\xi,c,d)\in\mathbb R^D$, $D=4L+31$,
where $c\in\mathbb R^{10}$ and $d\in\mathbb R^{24}$ are the components of the
curvature and of its gradient in the bases of Sect.~\ref{estimator} (there
$d$ is called $g$, a letter needed here for the metric). Let
$\sigma(\theta)\in\mathbb R^m$ be the model \eqref{model} for the pairs of
$\mathcal P$, with $\sigma_2$ as in \eqref{square} or with $\sigma_2=0$, and
$J(\theta)=\partial\sigma/\partial\theta$. We assume that the $m\times D$
matrix $J(\xi,0,0)$ has rank $D$ for every $\xi\in\mathcal K$.
\item[(A5) Weights.] The fit minimizes
$\chi^2(\theta)=N[y-\sigma(\theta)]^{\!\top}\widehat W[y-\sigma(\theta)]$
with a symmetric matrix $\widehat W$ that may depend on the data, and there
are constants $0<w_-\le w_+$ such that $w_-I\preceq\widehat W\preceq w_+I$
with probability at least $1-\varpi_N$, where $\varpi_N\to0$ as
$N\to\infty$. For the weight of Sect.~\ref{estimator},
$\widehat W=(N\widehat\Sigma)^{-1}$, this is proved in
Lemma~4 below.
\end{description}

Assumptions (A3) and (A4) are substantive, and (A4) is a condition on the
shape of the configuration only. By compactness there are $\tau_0>0$ and
$\mu>0$ such that $\eta(x_j-x_i,x_j-x_i)\le-\tau_0^2$ for
$(i,j)\in\mathcal P$ and the smallest singular value of $J(\xi,0,0)$ is at
least $\mu$, for all $\xi\in\mathcal K$; all constants below depend on
these numbers. The spatial axes of the frame $(e_a)$ depend on the
landmarks; since the orthonormal frames at $o$ with the given $e_0$ form a
compact set, all estimates below are uniform over them. The true parameter
is $\theta^*_T=(\xi,c^*,d^*)$, where $\xi$
is the true configuration and $c^*$, $d^*$ are the components of
$T^2C_{abcd}(o)$ and of $T^3\nabla_eC_{abcd}(o)$ in the frame $(e_a)$. The
measured clocks are $y_{ij}=-(n_{ij}/N)^{1/2}$, with $n_{ij}$ the number of
measurement events in $I(\ell_i,\ell_j)$, and the continuum clocks are
$\bar y_{ij}=-[V(\ell_i,\ell_j)/V(p,q)]^{1/2}$.

\subsection{Statement}
\noindent\textbf{Proposition 1.} \emph{Assume (A1)--(A5). There are
constants $T_1,r_1>0$ and $C_1,C_2<\infty$, which depend only on the
geometry of $(U,g)$, on $\mathcal K$, $\mathcal P$ and on $w_\pm$, such that
for all $0<T\le T_1$, all $N\ge2$ and every landmark configuration allowed
by (A3) the following holds with probability at least $1-\varpi_N-C_2/N$.}
\begin{enumerate}
\item[(i)] \emph{$\chi^2$ has exactly one critical point
$\hat\theta=(\hat\xi,\hat c,\hat d)$ in the closed ball of radius $r_1$
around $\theta^*_T$, and $\hat\theta$ is the unique minimizer of $\chi^2$
on this ball.}
\item[(ii)] \emph{$|\hat\theta-\theta^*_T|\le C_1\left(T^4+|y-\bar y|\right)$.}
\end{enumerate}
\emph{Moreover $\operatorname{E}|y-\bar y|^2\le m/N$. Consequently, for
every $t>0$, with probability at least $1-\varpi_N-C_2/N-m/t^2$,}
\begin{equation}
 |\hat\theta-\theta^*_T|\le C_1\left(T^4+tN^{-1/2}\right),\qquad
 \left|\hat c/T^2-c^*/T^2\right|\le C_1\left(T^2+tN^{-1/2}T^{-2}\right),
 \label{conv:rate}
\end{equation}
\emph{where $c^*/T^2$ are the components of the Weyl tensor $C_{abcd}(o)$.
The same holds, with other constants, in the tidal norm of
Sect.~\ref{setting}.}

\medskip
In short, $\hat\theta-\theta^*_T=O(T^4)+O_p(N^{-1/2})$ in the normalized
variables of the fit. The statement and its proof are the same with and
without $\sigma_2$, because $\sigma_2=O(|c|^2)$; in particular the
coefficient of the volume term in \eqref{square} is not used. The estimator
returns $\hat c$, an estimate of $T^2C(o)$. The second inequality of
\eqref{conv:rate} refers to the tensor in physical units, which requires
$T$, or equivalently bounds the relative error when $C(o)\neq0$.

\medskip
\noindent\textbf{Corollary 1.} \emph{Let the measurement events be a
Poisson sprinkling of density $\rho$ in $I(p,q)$, with lengths measured in
a fixed unit, and let (A1)--(A5) hold. If $T\to0$ and $\rho T^8\to\infty$,
then the solution $\hat\theta$ of Proposition~1 exists with probability
tending to one, and}
\begin{equation}
 \hat c/T^2-C(o)=O(T^2)+O_p\!\left(\rho^{-1/2}T^{-4}\right)\to0
 \quad\text{in probability}.
\end{equation}
\emph{For $T$ proportional to $\rho^{-1/12}$ the error is
$O_p(\rho^{-1/6})$.}

\subsection{Proofs}
\noindent\textbf{Lemma 1} (normal coordinates). \emph{Let $\mathcal F$ be a
compact set of orthonormal frames $(x,f)$ with $x\in U$. Let
$w\in\mathbb R^4$ be the normal coordinates based at $(x,f)$, and let
$s_{x,f}(w,w')$ and $\gamma_{x,f}(w)$ be the function $s$ and the matrix of
metric components in these coordinates. There are $r_0>0$ and $C_0$ such
that for all $(x,f)\in\mathcal F$ and $|w|,|w'|\le r_0$, with
$\Delta=w-w'$:}
\begin{enumerate}
\item[(a)] \emph{$s_{x,f}(w,w')=\eta(\Delta,\Delta)-\tfrac13R(w,w',w,w')
-\tfrac1{12}(w+w')^e\nabla_eR(w,w',w,w')+\varrho$ with
$|\varrho|\le C_0(|w|+|w'|)^6$, where $R$ and $\nabla R$ are the components
of the curvature and its gradient at $x$ in the frame $f$;}
\item[(b)] \emph{$s_{x,f}(w,w')-\eta(\Delta,\Delta)$ is bounded by
$C_0(|w|+|w'|)^4$ and its first derivatives by $C_0(|w|+|w'|)^3$;}
\item[(c)] $\bigl|\,|\det\gamma_{x,f}(w)|^{1/2}-1\,\bigr|\le C_0|w|^4$;
\item[(d)] $s_{x,f}(0,w)=\eta(w,w)$.
\end{enumerate}

\noindent\emph{Proof.} The maps $(x,f,w,w')\mapsto s_{x,f}(w,w')$ and
$(x,f,w)\mapsto\gamma_{x,f}(w)$ are smooth on an open set that contains
$\mathcal F\times\{0\}\times\{0\}$, hence on
$\mathcal F\times\{|w|,|w'|\le r_0\}$ for some $r_0>0$. Statement (d) is the
definition of $s$. For (a) apply Taylor's theorem in $(w,w')$ at the
origin: the derivatives of order six are bounded on this compact set, and
the Taylor polynomial of degree five is the expansion of the squared
geodesic interval in normal coordinates, Sect.~9 of Ref.~\cite{brewin}
(Eqs.~(11.21a)--(11.21c) in arXiv:0903.2087v3). That polynomial has no terms
of degree three, which gives (b) in the same way. For (c),
$\gamma_{ab}(w)=\eta_{ab}-\tfrac13R_{acbd}w^cw^d
-\tfrac16\nabla_eR_{acbd}w^cw^dw^e+O(|w|^4)$ \cite{brewin}, and the traces
of the second and third terms with $\eta^{ab}$ contain only the Ricci tensor
and its gradient, which vanish; hence $\det\gamma=-1+O(|w|^4)$, with a
remainder bounded as before. \hfill$\square$

\medskip
\noindent\textbf{Lemma 2} (volume of a small vacuum interval). \emph{Let
$\mathcal T$ be a compact set of future-directed unit timelike vectors
$(x,u)$ with $x\in U$. There are $\tau_1>0$ and $C_V$ such that, whenever
$(x,u)\in\mathcal T$, $0<\tau\le\tau_1$, $x'=\exp_x(\tau u)$ and
$J^+(x)\cap J^-(x')\subset U$,}
\begin{equation}
 \left|V(x,x')-\tfrac{\pi}{24}\tau^4\right|\le C_V\tau^8 .
 \label{conv:volume}
\end{equation}

\noindent\emph{Proof.} Complete $u$ to an orthonormal frame $f$ with
$f_0=u$. The frames so obtained form a compact set, to which Lemma~1
applies. We use the normal coordinates $w$ based at $(x,f)$, in which $x'$
is $\tau\hat e$ with $\hat e=(1,0,0,0)$, write $s=s_{x,f}$, and put
$\Xi(w)=s(0,w)-s(w,\tau\hat e)$. We show that the half
$H_-=\{z\in J^+(x)\cap J^-(x'):\Xi\ge0\}$ of the interval has volume
$\tfrac{\pi}{48}\tau^4+O(\tau^8)$. The other half, where $\Xi\le0$, is
treated in the same way from $x'$ with the time orientation reversed, and
the common boundary $\Xi=0$ has measure zero. All $O(\cdot)$ below are
uniform over $\mathcal T$ and over the directions $\nu$.

(1) \emph{Localization.} By (A1), every causal vector $v$ at a point of $U$
has $|\vec v|\le\tfrac54|v^0|$ in the coordinates of (A1), with $v^0>0$ if
it is future-directed. Hence $x^0$ increases along future-directed causal
curves in $U$, and the coordinate differences satisfy
$|z-x|\le\tfrac94|x'-x|$ for every $z\in J^+(x)\cap J^-(x')$. Since the
change to the coordinates $w$ is smooth, uniformly over the frames, there
is $C\ge1$ such that $|w|\le C\tau$ for all these $z$.

(2) \emph{Rays.} In a convex normal neighbourhood, $z\in J^+(x)$ if and
only if $\exp_x^{-1}z$ is zero or future-directed causal
\cite[Lemma~14.2]{oneill}, that is, if and only if $w=t\nu$ with $t\ge0$
and $\nu\in\mathcal N=\{(1,\vec\nu):|\vec\nu|\le1\}$. For $\nu\in\mathcal N$
let $\phi_\nu(t)=\Xi(t\nu)$. By Lemma~1(d),
$\phi_\nu(t)=t^2\eta(\nu,\nu)-s(t\nu,\tau\hat e)$, and
$\eta(t\nu-\tau\hat e,t\nu-\tau\hat e)=t^2\eta(\nu,\nu)+2\tau t-\tau^2$.
By Lemma~1(b), for $0\le t\le C\tau$,
\begin{equation}
 \phi_\nu(t)=\tau^2-2\tau t+O(\tau^4),\qquad
 \phi_\nu'(t)=-2\tau+O(\tau^3).
\end{equation}
For $\tau\le\tau_1$ small, $\phi_\nu$ is therefore strictly decreasing on
$[0,C\tau]$, with $\phi_\nu(0)=\tau^2$ and $\phi_\nu(\tau)<0$. It has a
single zero $t_*(\nu)\in(0,\tau)$, which depends continuously on $\nu$.

(3) \emph{The half-interval is
$\{t\nu:\nu\in\mathcal N,\ 0\le t\le t_*(\nu)\}$.} If $z\in H_-$, then
$w=t\nu$ with $t\le C\tau$ by (1) and (2), and $\phi_\nu(t)\ge0$ gives
$t\le t_*(\nu)$. Conversely, let $0\le t\le t_*(\nu)$. On the segment from
$0$ to $t\nu$ we have $\Xi\ge0$, that is, $s(\cdot,x')\le s(x,\cdot)\le0$,
so the segment lies in the set of events causally related to $x'$, and it
does not contain $x'$, where $\Xi=-\tau^2$. The sets
$J^-(x')\setminus\{x'\}$ and $J^+(x')\setminus\{x'\}$ are disjoint and
closed in $U\setminus\{x'\}$, and the segment starts at $x\in J^-(x')$;
being connected, it lies in $J^-(x')$. Hence $t\nu\in H_-$.

(4) \emph{Volume.} With $w=(t,t\vec\nu)$ we have
$\mathrm d^4w=t^3\,\mathrm dt\,\mathrm d^3\vec\nu$, and Lemma~1(c) gives
\begin{equation}
 \operatorname{Vol}(H_-)=\int_{|\vec\nu|\le1}\int_0^{t_*(\nu)}
 \left[1+O(\tau^4)\right]t^3\,\mathrm dt\,\mathrm d^3\vec\nu
 =\frac14\int_{|\vec\nu|\le1}t_*(\nu)^4\,\mathrm d^3\vec\nu+O(\tau^8).
 \label{conv:half}
\end{equation}

(5) \emph{Expansion of $t_*$.} Insert Lemma~1(a) with $w=t\nu$ and
$w'=\tau\hat e$. Since
$R(t\nu,\tau\hat e,t\nu,\tau\hat e)=t^2\tau^2E_{ij}\nu^i\nu^j$ with
$E_{ij}=R_{0i0j}$,
\begin{equation}
 \phi_\nu(t)=\tau^2-2\tau t+\tfrac13t^2\tau^2E_{ij}\nu^i\nu^j
 +\tfrac1{12}t^2\tau^2\left[(t+\tau)\nabla_0E_{ij}
 +t\,\nu^k\nabla_kE_{ij}\right]\nu^i\nu^j+O(\tau^6).
\end{equation}
The zero satisfies $t_*=\tfrac\tau2+O(\tau^3)$, and then
\begin{equation}
 t_*(\nu)=\frac\tau2\left[1+\frac{\tau^2}{12}E_{ij}\nu^i\nu^j
 +\frac{\tau^3}{96}\left(3\nabla_0E_{ij}+\nu^k\nabla_kE_{ij}\right)
 \nu^i\nu^j+O(\tau^4)\right].
\end{equation}
The integral of $\nu^i\nu^j$ over the unit ball is
$\tfrac{4\pi}{15}\delta^{ij}$, and that of $\nu^i\nu^j\nu^k$ vanishes. In
vacuum $E_{ii}=R_{00}=0$ and $\nabla_0E_{ii}=\nabla_0R_{00}=0$, so the terms
of relative order $\tau^2$ and $\tau^3$ integrate to zero, and
\eqref{conv:half} gives
$\operatorname{Vol}(H_-)=\tfrac14\cdot\tfrac{4\pi}3(\tau/2)^4+O(\tau^8)
=\tfrac{\pi}{48}\tau^4+O(\tau^8)$. \hfill$\square$

\medskip
Lemma~2 uses the field equations only through $R_{ab}=0$, and it gives only
the order of the correction. As a check of the method, without that
condition the same computation, with the volume element
$1-\tfrac16R_{ab}w^aw^b$, gives
$V=\tfrac{\pi}{24}\tau^4[1-\tfrac1{180}R^a{}_a\tau^2+\tfrac1{30}R_{00}\tau^2
+O(\tau^3)]$, which is the formula of Myrheim, Eqs.~(7), (14) and (23) of
Ref.~\cite{gibbons}. In vacuum the first correction is of relative order
$\tau^4$, in agreement with Eq.~(79) of Ref.~\cite{wang}; its coefficient, on
which \eqref{square} and Ref.~\cite{wang} differ, plays no role here.

\medskip
\noindent\textbf{Lemma 3} (clocks). \emph{Assume (A1)--(A3). (a) There are
$T_1>0$ and $C_b$ such that $|\bar y_{ij}-\sigma_{ij}(\theta^*_T)|\le
C_bT^4$ for all $T\le T_1$, $\xi\in\mathcal K$ and $(i,j)\in\mathcal P$.
(b) $\operatorname{E}|y-\bar y|^2\le m/N$.}

\noindent\emph{Proof.} (a) Lemma~1(a) at $(o,e)$ with $w=Tx_i$ and
$w'=Tx_j$, and $R_{abcd}=C_{abcd}$ in vacuum, give
\begin{equation}
 \frac{s(\ell_i,\ell_j)}{T^2}=\eta(\Delta,\Delta)
 -\frac{T^2}3C(x_i,x_j,x_i,x_j)
 -\frac{T^3}{12}(x_i+x_j)^e\nabla_eC(x_i,x_j,x_i,x_j)+O(T^4),
\end{equation}
with $\Delta=x_i-x_j$ and the tensors at $o$. The right-hand side is the
model \eqref{model} at $\theta^*_T$ without $\sigma_2$, up to $O(T^4)$, and
$\sigma_2(\theta^*_T)=O(|c^*|^2)=O(T^4)$. Since
$\eta(\Delta,\Delta)\le-\tau_0^2$, for small $T$ we have
$\tau_{ij}^2=-s(\ell_i,\ell_j)\ge\tau_0^2T^2/2$ and $\tau_{ij}=O(T)$. In
the coordinates of (A1) the vector $\exp_{\ell_i}^{-1}\ell_j$ has components
$T(x_j-x_i)+O(T^2)$, so for small $T$ it is future-directed and $\ell_j$
lies in the timelike future of $\ell_i$. The unit tangent at $\ell_i$ of the
geodesic to $\ell_j$ then has bounded components, and these vectors lie in a
compact set $\mathcal T$.
Moreover $J^+(\ell_i)\cap J^-(\ell_j)\subset I(p,q)\subset U$. Lemma~2
therefore gives $V(\ell_i,\ell_j)=\tfrac{\pi}{24}\tau_{ij}^4[1+O(T^4)]$, and
likewise $V(p,q)=\tfrac{\pi}{24}T^4[1+O(T^4)]$. Hence
$\bar y_{ij}=-(\tau_{ij}^2/T^2)[1+O(T^4)]=s(\ell_i,\ell_j)/T^2+O(T^4)$.

(b) Given the landmarks, $n_{ij}$ is binomial with parameters $N$ and
$\pi_{ij}=V(\ell_i,\ell_j)/V(p,q)=\bar y_{ij}^2$. Since
$|\sqrt a-\sqrt b|\le|a-b|/\sqrt b$,
\begin{equation}
 |y_{ij}-\bar y_{ij}|\le\frac{|n_{ij}/N-\pi_{ij}|}{\sqrt{\pi_{ij}}},
 \qquad\operatorname{E}|y_{ij}-\bar y_{ij}|^2\le\frac{1-\pi_{ij}}N .
\end{equation}
Summing over the $m$ pairs gives (b). The clocks are strongly correlated,
see \eqref{cov}, but no independence is used. \hfill$\square$

\medskip
\noindent\textbf{Lemma 4} (weights). \emph{Let $N\widehat\Sigma=\widehat
D\,\widehat\Gamma\,\widehat D$ be the covariance of Sect.~\ref{estimator},
where $\widehat\Gamma$ is the empirical covariance matrix of the membership
indicators of the $m$ intervals and
$\widehat D=\operatorname{diag}(\tfrac12\hat\pi_{ij}^{-1/2})$ with
$\hat\pi_{ij}=n_{ij}/N$. Under (A1)--(A3) there are $T_2>0$,
$0<w_-\le w_+$ and $C_W$ such that $\widehat W=(N\widehat\Sigma)^{-1}$
satisfies (A5) with $\varpi_N=C_W/N$ for all $T\le T_2$.}

\noindent\emph{Proof.} Let $\Gamma_T(\xi)$ be the covariance matrix of the
indicators under the sampling law, and $\Gamma_0(\xi)$ the same matrix for
the flat metric $\eta$ and the points $x_i$.

(1) \emph{$\Gamma_0(\xi)$ is positive definite.} Otherwise a combination
$\sum\lambda_{ij}1_{ij}$ of the indicators with $\lambda\neq0$ is constant
almost everywhere in the flat interval. Fix a pair $(i,j)$. The edge of its
interval, where the future cone of $x_i$ meets the past cone of $x_j$, is a
round two-sphere, and the only null cones that contain it are those of
$x_i$ and $x_j$; the null cone of any other point meets it in the zero set
of a quadratic polynomial that does not vanish identically on the sphere,
a closed set without interior points. Choose $z$ on the edge and on none of
the other cones, and a ball $B$ around $z$ that meets no null cone of a
point $x_k$ other than the future cone of $x_i$ and the past cone of $x_j$.
In $B$, every indicator other than $1_{ij}$ is a constant, or a constant
times the indicator of $J^+(x_i)$ or of $J^-(x_j)$, while $1_{ij}$ is the
product of these two. The two cones cross transversally at $z$. If
$1\le i,j\le L$, all four combinations of the two indicators occur on sets
of positive measure in $B$, and constancy of the sum forces
$\lambda_{ij}=0$. If $i=0$, only $B\cap J^+(p)$ carries probability, and no
interval $I(x_k,x_j)$ with $k\neq0$ meets $B$, because
$z\in\partial J^+(p)$ whereas $x_k$ lies in the open interval; the two
values of the indicator of $J^-(x_j)$ again force $\lambda_{0j}=0$. The
case $j=L+1$ is symmetric.

(2) In the coordinates $x/T$ the metric has components $g_{ab}(Tx)$. These,
the inverse exponential map and the volume density converge to their flat
values as $T\to0$, uniformly on compact sets. The indicator of each
interval is therefore continuous in $(T,\xi)$ and in the spatial frame at
$o$ at every point off the boundary of the limiting interval, and by
dominated convergence $\Gamma_T(\xi)$ is continuous in these variables on
$[0,T_0]\times\mathcal K$ times the compact set of frames, with the value
$\Gamma_0(\xi)$ at $T=0$. With (1) and compactness,
$\Gamma_T(\xi)\succeq4\gamma I$ for some $\gamma>0$, all $\xi\in\mathcal K$,
all frames and all $T\le T_2$. In the same way $\pi_{ij}\ge\pi_->0$.

(3) $\widehat\Gamma$ and $\hat\pi$ are built from averages of bounded
variables, so $\operatorname{E}\|\widehat\Gamma-\Gamma_T\|^2\le Cm^2/N$ and
$\operatorname{E}|\hat\pi-\pi|^2\le m/N$. By Chebyshev's inequality the
event that $\|\widehat\Gamma-\Gamma_T\|\le2\gamma$ and
$\hat\pi_{ij}\ge\pi_-/2$ for all pairs has probability at least $1-C_W/N$.
On this event the diagonal entries of $\widehat D$ lie between $\tfrac12$
and $(2\pi_-)^{-1/2}$, $\widehat\Gamma\succeq2\gamma I$ and
$\operatorname{tr}\widehat\Gamma\le m/2$, hence
$\tfrac\gamma2I\preceq N\widehat\Sigma\preceq\tfrac{m}{4\pi_-}I$.
\hfill$\square$

\medskip
\noindent\textbf{Lemma 5} (local minimizer). \emph{Let $F$ be a twice
continuously differentiable map from the closed ball of radius $r_0$ around
$\theta^*\in\mathbb R^D$ to $\mathbb R^m$, with $\|DF\|\le M_1$ and
$\|D^2F\|\le M_2$. Let $W$ be symmetric with $0\preceq W\preceq w_+I$ and
$DF(\theta^*)^{\!\top}W\,DF(\theta^*)\succeq\lambda I$, $\lambda>0$. For
$y\in\mathbb R^m$ let $Q(\theta)=\tfrac12[y-F(\theta)]^{\!\top}
W[y-F(\theta)]$ and $\epsilon=|y-F(\theta^*)|$, and put}
\begin{equation}
 r_1=\min\left\{r_0,\frac{\lambda}{12M_1M_2w_+}\right\},\qquad
 \epsilon_1=\min\left\{\frac{\lambda}{4M_2w_+},
 \frac{\lambda r_1}{8M_1w_+}\right\}.
\end{equation}
\emph{If $\epsilon\le\epsilon_1$, then $Q$ has exactly one critical point
$\hat\theta$ in the closed ball of radius $r_1$ around $\theta^*$, it is the
unique minimizer of $Q$ on this ball, and
$|\hat\theta-\theta^*|\le4M_1w_+\epsilon/\lambda$.}

\noindent\emph{Proof.} Let $|\theta-\theta^*|\le r_1$, let $v$ be a unit
vector and $\rho=y-F(\theta)$. Then
\begin{equation}
 v^{\!\top}D^2Q(\theta)\,v=[DF(\theta)v]^{\!\top}W[DF(\theta)v]
 -\rho^{\!\top}W\,D^2F(\theta)[v,v].
\end{equation}
Here $|\rho|\le\epsilon+M_1r_1$ and $\|DF(\theta)-DF(\theta^*)\|\le M_2r_1$,
so that
$v^{\!\top}D^2Q(\theta)v\ge\lambda-2M_1M_2w_+r_1-M_2w_+(\epsilon+M_1r_1)
\ge\lambda/2$. Thus $Q$ is strongly convex on the ball. It has a unique
minimizer $\hat\theta$ there and at most one critical point, and
$Q(\theta)\ge Q(\theta^*)+DQ(\theta^*)(\theta-\theta^*)
+\tfrac\lambda4|\theta-\theta^*|^2$. Since
$|DQ(\theta^*)|=|DF(\theta^*)^{\!\top}W[y-F(\theta^*)]|\le M_1w_+\epsilon$,
we get $Q(\theta)>Q(\theta^*)$ whenever
$|\theta-\theta^*|>4M_1w_+\epsilon/\lambda$. Hence
$|\hat\theta-\theta^*|\le4M_1w_+\epsilon/\lambda\le r_1/2$, so $\hat\theta$
is an interior point and therefore a critical point. \hfill$\square$

\medskip
\noindent\emph{Proof of Proposition 1.} The map $\sigma$ is a polynomial in
$\theta$. Let $M_1$ and $M_2$ bound its first and second derivatives on the
set of all $\theta$ at distance at most one from
$\{(\xi,c,d):\xi\in\mathcal K,\ |c|,|d|\le1\}$; for small $T$ this set
contains the ball of radius $r_0=1$ around $\theta^*_T$. Because $\sigma_2$
is quadratic in $c$, the matrix $J(\xi,0,0)$ does not depend on it, and
$\|J(\theta^*_T)-J(\xi,0,0)\|\le M_2(|c^*|+|d^*|)=O(T^2)$. By (A4) the
smallest singular value of $J(\theta^*_T)$ is at least $\mu/2$ for
$T\le T_1$, so that on the event of (A5)
$J(\theta^*_T)^{\!\top}\widehat WJ(\theta^*_T)\succeq\lambda I$ with
$\lambda=w_-\mu^2/4$. Apply Lemma~5 on this event with $F=\sigma$,
$W=\widehat W$ and $\theta^*=\theta^*_T$; the functions $\chi^2$ and $Q$
differ by the factor $2N$. By Lemma~3(a),
$\epsilon=|y-\sigma(\theta^*_T)|\le|y-\bar y|+\sqrt m\,C_bT^4$. Reduce $T_1$
so that $\sqrt m\,C_bT_1^4\le\epsilon_1/2$. By Lemma~3(b) and Chebyshev's
inequality, $|y-\bar y|\le\epsilon_1/2$ except on an event of probability at
most $C_2/N$ with $C_2=4m/\epsilon_1^2$. Lemma~5 then gives (i), and (ii)
with $C_1=4M_1w_+\max(1,\sqrt m\,C_b)/\lambda$. Finally
$|y-\bar y|\le tN^{-1/2}$ except with probability $m/t^2$, which gives
\eqref{conv:rate}, the second inequality because
$|\hat c-c^*|\le|\hat\theta-\theta^*_T|$. \hfill$\square$

\medskip
\noindent\emph{Proof of Corollary 1.} Given their number $N$, the events of
a Poisson sprinkling in $I(p,q)$ satisfy (A2), and $N$ is Poisson with mean
$\rho V(p,q)=\tfrac{\pi}{24}\rho T^4[1+O(T^4)]$ by Lemma~2. If
$\rho T^4\to\infty$, then $N\ge\rho T^4/20$ with probability tending to
one, and on this event $N^{-1/2}T^{-2}\le\sqrt{20}\,\rho^{-1/2}T^{-4}$ and
$\varpi_N+C_2/N\to0$. Insert this into \eqref{conv:rate}. The two terms
balance for $T^6$ proportional to $\rho^{-1/2}$, and then
$\rho T^8\to\infty$. \hfill$\square$

\subsection{What is not proved}\label{app:open}
\emph{The selection rule.} Assumptions (A3) and (A4) are not shown to hold
for the landmarks produced by the rule of Sect.~\ref{estimator}, which are
events of a selection sample chosen by counts, with a clock set chosen by
counts. Assumption (A4) is not vacuous: for an explicit rational
configuration with five landmarks in each of three time bands, and with
the clocks between different bands and to the tips, exact modular
arithmetic gives $J(\xi,0,0)$ the full rank $D=91$. Since full rank is the
non-vanishing of a polynomial, (A4) then holds for almost every
configuration with these clocks (see Corollary~2). This identifies the
parameters; it does not bound $\mu$, on which the constants depend. In
that configuration the singular values of $J(\xi,0,0)$ range over eight
orders of magnitude.

\emph{The optimizer.} Proposition~1 concerns the critical point near
$\theta^*_T$. It does not show that the Levenberg--Marquardt iteration from
the four starts of Sect.~\ref{estimator} converges to this point, and it
does not exclude other local minima. The reflections of
Sect.~\ref{setting} map $\hat\theta$ to further minimizers with the same
$\chi^2$; these are the mirror images accounted for there. No global
identifiability is proved.

\emph{The coverage.} Nothing is proved about the ball \eqref{ball}. For
$T$ proportional to $\rho^{-1/12}$ the bias and the noise in
\eqref{conv:rate} are of the same order, so a ball built from the noise
covariance alone cannot be expected to cover, and the allowance $\delta$
remains an assumption. In the regime $N^{1/2}T^4\to0$, where the noise
dominates, one may expect $\hat\theta$ to be asymptotically normal with the
covariance used in Sect.~\ref{estimator}, but we have not proved this, and
no asymptotic statement bears on the coverage at the values of $N$ and $T$
of the tests.

\emph{The constants.} $C_1$ and $T_1$ depend on $\mu$ and on bounds for the
derivatives of the metric on $U$ up to the order that controls the sixth
derivatives of $s$. They are not computed, so Proposition~1 gives no bias
bound at a given window, such as the allowance of 2\% in \eqref{ball}. It
is not uniform over spacetimes without a bound on these derivatives, in
agreement with the obstruction of Sect.~\ref{limits}.

\emph{Further restrictions.} The tips are given events and are not chosen
from the causal set. Vacuum is essential: Ricci terms change the clocks
\eqref{clock} at relative order $T^2$, the order of the signal. The
implementation discards eigenvalues of the clock correlation matrix below
$10^{-8}$ times the largest. Lemma~4 covers this only when no eigenvalue is
discarded; otherwise the lower bound in (A5) has to be replaced by a lower
bound on $J^{\!\top}\widehat WJ$, which is all that the proof uses. The rate
is an upper bound for this estimator under these assumptions, and we claim
no optimality.

\subsection{An idealized random design}\label{app:random}
The conditioning assumption can be removed, in probability, for a
simplified design in which the landmarks are placed at random with a
bounded density. This is not the order-based rule of
Sect.~\ref{estimator}.

\medskip
\noindent\textbf{Corollary 2.} \emph{Let $B_-,B_0,B_+$ be compact subsets
of the open flat interval $\{|x^0|+|\vec x|<\tfrac12\}$ such that $x'-x$ is
future-directed timelike with respect to $\eta$ whenever
$x\in B_-,x'\in B_0$, or $x\in B_0,x'\in B_+$, or $x\in B_-,x'\in B_+$.
Assume (A1) and (A2), and let the landmarks be random, with at least five
in each of the three sets $\exp_o(TB_-)$, $\exp_o(TB_0)$, $\exp_o(TB_+)$,
their numbers being fixed. Assume that their normal coordinates in units
of $T$, with respect to a fixed frame at $o$, have a joint density bounded
by a constant independent of $T$. Let the clocks be all pairs of landmarks
from different sets and all pairs of a landmark and a tip, and let
$\widehat W$ be the weight of Lemma~4 or the identity. Then for every
$\varepsilon>0$ there are $r,C,T_\varepsilon,N_\varepsilon$ such that, for
$T\le T_\varepsilon$ and $N\ge N_\varepsilon$, statements (i) and (ii) of
Proposition~1 hold with $r_1=r$ and $C_1=C$ with probability at least
$1-\varepsilon$. In particular
$\hat\theta-\theta^*_T=O_p(T^4+N^{-1/2})$.}

\noindent\emph{Proof.} Let $\tilde x=(\tilde x_1,\dots,\tilde x_L)$ be the
coordinates in the fixed frame; they range over a compact product
$\mathcal S$ of the three sets. Let $\tilde J(\tilde x)$ be the Jacobian of
the model with respect to all $4L$ coordinates and $(c,d)$, at $c=d=0$. Its
entries are polynomials in $\tilde x$. The model is invariant under a
common spatial rotation of the points and of the tensors, so the three
infinitesimal rotations of the points lie in the kernel of $\tilde J$, and
$\tilde J$ has the same rank at $\tilde x$ and at the rotated configuration
$x$ that satisfies the gauge of (A3). When $x_1^1\neq0$ and $x_2^2\neq0$,
the three gauge conditions define a complement of the infinitesimal
rotations, and therefore $J(\xi,0,0)$ has rank $D$ if and only if
$\tilde J(\tilde x)$ has rank $4L+31$. The configurations of smaller rank
are the common zeros of the minors of order $4L+31$. These polynomials do
not all vanish identically: for five landmarks in each set this is the
exact computation quoted in Sect.~\ref{app:open}, and a further landmark
adds four columns, which are nonzero only in the rows of its own clocks,
where they equal $\pm2\eta(x_{\rm new}-x_j)$ and have rank four for a
suitable position, so that the rank increases by four. Hence the
configurations of smaller rank, together with those in which two points
coincide or the spatial parts of $\tilde x_1$ and $\tilde x_2$ are linearly
dependent, form a set $\mathcal Z$ of Lebesgue measure zero.

For $\mu>0$ let $\mathcal G_\mu\subset\mathcal S$ be the set on which the
singular value of order $4L+31$ of $\tilde J$, the distances between the
points, $|\vec{\tilde x}_1|$ and the area spanned by $\vec{\tilde x}_1$
and $\vec{\tilde x}_2$ are all at least $\mu$. These sets are compact,
they increase to $\mathcal S\setminus\mathcal Z$ as $\mu\to0$, and the
gauge rotation is continuous on them. Their images $\mathcal K_\mu$ in the
gauge of (A3) are compact and satisfy (A3) and (A4), and for small $T$ the
landmarks lie in $I(p,q)$ by Lemma~1. Since the density is bounded, the
probability that $\tilde x\notin\mathcal G_\mu$ is at most a constant times
the Lebesgue measure of $\mathcal S\setminus\mathcal G_\mu$, which tends to
zero as $\mu\to0$, uniformly in $T$. Given $\varepsilon$, choose $\mu$ such
that this probability is at most $\varepsilon/2$, and apply Proposition~1
and Lemma~4 with $\mathcal K=\mathcal K_\mu$, conditionally on the
landmarks. \hfill$\square$

\medskip
The constants of Corollary~2 depend on $\varepsilon$ through $\mu$ and are
not controlled; the remarks of Sect.~\ref{app:open} on the optimizer, the
coverage and the constants apply unchanged.

\subsection{Relation to the theorem of Madsen}
Theorem~4.9 and Corollary~4.10 of Ref.~\cite{madsen} hold for a smooth
globally hyperbolic spacetime of any dimension $d\ge2$, without field
equations. The input is a Poisson sprinkling of known density in a fixed
compact causal diamond. Anchors are selected from the sample by tests on
counts, and the existence of accepted anchors is proved. The output is the
full Riemann tensor at every sample point that passes the tests, in a
reconstructed frame, and the bound
$C[\log(e+\rho)/\rho]^{1/[2(3d+4)]}$ holds simultaneously at all these
points with probability at least $1-\rho^{-n}$ for every $n$, with a
constant that depends on $n$, and almost surely for large $\rho$.

Proposition~1 is weaker in each of these respects. It is restricted to
$d=4$ and to vacuum, the tips are given, the conditioning of the landmarks
and the choice of the local solution are assumed, the tensor is estimated
at one point, and the convergence is in probability. Under these
assumptions the exponent is $1/6$ instead of $1/32$. The two exponents do
not refer to the same problem, and the comparison does not show that one
construction is more efficient than the other. In particular, the rate
$\rho^{-1/6}$ relies on vacuum, which makes the volume clocks accurate to
relative order $T^4$ (Lemma~2); without it the clocks are accurate only to
relative order $T^2$ and the present estimator is not consistent.

\subsection{Numerical check of the orders}
The script \texttt{check\_convergence\_orders.py} in the repository checks
the orders used above on exact rotating plane waves
(Sect.~\ref{tests}), for a fixed configuration of 48 landmarks in windows
scaled by $1$, $\tfrac12$, $\tfrac14$ and $\tfrac18$ and with exact
interval volumes in place of counts. Over the three successive halvings of
the window, the residual of the fifth-degree expansion of Lemma~1(a)
decreases by factors of $15.1$, $15.8$ and $15.9$, and the relative volume
correction of Lemma~2 by $15.8$, $15.9$ and $16.0$, to be compared with
$16$. The relative error of the fitted tensor decreases by $4.3$, $4.1$
and $4.0$ without $\sigma_2$ and by $3.8$, $4.0$ and $4.0$ with it, to be
compared with $4$, and the error $|\hat\theta-\theta^*_T|$ of all
parameters by $22.8$, $18.0$ and $16.6$ without $\sigma_2$ and by $26.9$,
$17.4$ and $16.6$ with it, to be compared with $16$. The same script
verifies the statement on Myrheim's formula made after Lemma~2 and
performs the exact rank computation
quoted in Sect.~\ref{app:open}. This checks the orders in one family of
spacetimes and one configuration; it does not test the assumptions.
